% REPO_EXPLAINED.tex — a technical explanation of the compliant-wheel repository.
% Compiles with pdflatex (twice, for the table of contents and cross-references):
%     pdflatex REPO_EXPLAINED.tex && pdflatex REPO_EXPLAINED.tex
% or simply:  latexmk -pdf REPO_EXPLAINED.tex

\documentclass[11pt,a4paper]{article}

\usepackage[T1]{fontenc}
\usepackage[utf8]{inputenc}
\usepackage{lmodern}
\usepackage[margin=0.95in]{geometry}
\usepackage{amsmath}
\usepackage{amssymb}
\usepackage{booktabs}
\usepackage{longtable}
\usepackage{tabularx}
\usepackage{array}
\usepackage{xcolor}
\usepackage{listings}
\usepackage{microtype}
\usepackage{enumitem}
\usepackage[hidelinks,bookmarksnumbered]{hyperref}

% ---------------------------------------------------------------------------
% Conventions
% ---------------------------------------------------------------------------

% Inline code.  Underscores are escaped explicitly throughout (\_), and \_ is
% locally redefined so that a long identifier may break across lines after an
% underscore -- the same courtesy \url grants a long URL.  Without this the
% glossary's narrow middle column overflows on names like
% \code{wheel\_step\_export.\_embed}.
\DeclareRobustCommand{\code}[1]{%
  \texttt{\small\def\_{\textunderscore\discretionary{}{}{}}#1}}

% A file-and-line citation, e.g. \at{wheel\_wheel.py}{2209}
\newcommand{\at}[2]{\code{#1:#2}}

\definecolor{codebg}{gray}{0.96}
\definecolor{codeframe}{gray}{0.80}
\definecolor{codecmt}{rgb}{0.25,0.45,0.25}
\definecolor{codekw}{rgb}{0.20,0.20,0.60}

\lstdefinestyle{plain}{
  basicstyle=\ttfamily\footnotesize,
  backgroundcolor=\color{codebg},
  frame=single,
  rulecolor=\color{codeframe},
  framesep=5pt,
  xleftmargin=6pt,
  xrightmargin=6pt,
  keepspaces=true,
  columns=fullflexible,
  breaklines=false,
  showstringspaces=false,
  captionpos=b,
  aboveskip=1.0em,
  belowskip=1.0em,
}

\lstdefinestyle{py}{
  style=plain,
  language=Python,
  keywordstyle=\color{codekw}\bfseries,
  commentstyle=\color{codecmt}\itshape,
  stringstyle=\color{codecmt},
}

\lstset{style=plain}

% Long unbreakable identifiers in narrow columns: let TeX stretch the line rather
% than run it off the page.
\setlength{\emergencystretch}{2.5em}

% Tighter lists.
\setlist[itemize]{itemsep=0.25em,topsep=0.4em}
\setlist[enumerate]{itemsep=0.25em,topsep=0.4em}

% Glossary tables: three columns, the middle one doing the work.
\newenvironment{termtable}[1]
  {\small\begin{longtable}{@{}>{\raggedright\arraybackslash}p{0.19\linewidth}
                             >{\raggedright\arraybackslash}p{0.54\linewidth}
                             >{\raggedright\arraybackslash}p{0.21\linewidth}@{}}
   \caption{#1}\\
   \toprule
   \textbf{Term} & \textbf{Meaning} & \textbf{Where} \\
   \midrule
   \endfirsthead
   \toprule
   \textbf{Term} & \textbf{Meaning} & \textbf{Where} \\
   \midrule
   \endhead
   \midrule
   \multicolumn{3}{r@{}}{\emph{continued on the next page}}\\
   \endfoot
   \bottomrule
   \endlastfoot}
  {\end{longtable}}

\title{\textbf{The Compliant Wheel Repository}\\[0.35em]
       \large What it does, and the mathematics it runs on}
\author{A guided tour of \code{github/wheel}}
\date{Written 2026-08-29 against the \code{feature} branch}

\begin{document}
\maketitle

\begin{abstract}
\noindent
A tour of this repository for a reader who is technically fluent but new both to the
tree and to structural-optimization mathematics. It explains what the software
designs, the specific mathematics each stage implements, the project's private
vocabulary, the module architecture, an end-to-end trace of one optimizer step, and
the reasoning behind the algorithmically interesting choices. Every claim is traceable
to a file and a line; where the code and the plan files disagree, the code is followed.
Line references are against the \code{feature} branch as of 2026-08-29, with
\code{PLAN.md} sections referenced up to \S92.
\end{abstract}

\tableofcontents
\clearpage

% ===========================================================================
\section{Overview}
\label{sec:overview}
% ===========================================================================

\subsection{The thing being designed}

This repository designs \emph{one physical object}: a compliant (deliberately springy)
wheel for a small UAV, printed in PLA on a fused-filament (FFF) printer. It is a flat,
planar part extruded to a constant face width---you could cut it out of a 22.4\,mm
thick plate. Its geometry is fixed by a handful of hard requirements
(\at{src/wheel\_fea.py}{107--165}):

\begin{table}[h]
\centering
\small
\begin{tabular}{@{}lll@{}}
\toprule
\textbf{Quantity} & \textbf{Value} & \textbf{Where} \\
\midrule
outer diameter    & $\varnothing$100\,mm (\code{RIM\_OUTER\_RADIUS\_MM = 50.0}) & \at{wheel\_wheel.py}{173} \\
hub bore radius   & 12.7\,mm (a $\tfrac12''$ shaft)                             & \at{wheel\_fea.py}{110} \\
spokes            & 12, twelve-fold rotationally symmetric                      & \at{wheel\_fea.py}{109} \\
face width        & 22.4\,mm                                                    & \at{wheel\_fea.py}{145} \\
material          & PLA, $E = 2300$\,MPa, $\nu = 0.35$                          & \at{wheel\_fea.py}{148} \\
service load      & 15\,lbf $=$ 66.72\,N                                        & \at{wheel\_fea.py}{152--153} \\
\textbf{target axle drop} & \textbf{2.0\,mm}                                    & \at{wheel\_fea.py}{158} \\
allowable stress  & 25\,MPa (50 ult.\ $\times$ 0.80 FFF $\div$ 1.6 SF)          & \at{wheel\_fea.py}{149--151} \\
\bottomrule
\end{tabular}
\caption{The fixed requirements. Everything else is optimized.}
\end{table}

The interesting entry is the target deflection. This is not a stiff wheel that must not
bend; it is a \emph{compliant mechanism}---the 2.0\,mm of travel under load \textbf{is
the feature}, a suspension made of geometry instead of springs. That is why the
deflection objective is two-sided: being too stiff is penalised exactly as hard as
being too soft (\at{wheel\_fea.py}{748--752}, \at{wheel\_objective.py}{1152--1158}).

Between the hub and the rim run twelve identical curved spokes. Their shape is what
gets optimized.

\subsection{What the software actually does}

The repository is a \emph{design-optimization pipeline with a verification apparatus
bolted to every joint}. It has three stages plus an exporter.

\begin{lstlisting}
   Stage 1/2  GENETIC ALGORITHM over an analytical beam model
              wheel_fea.py + pygad          ~minutes         14-gene genome
                    |
                    |   best_solution.json / stage2_elites.json
                    v
   Stage 3    GRADIENT DESCENT over a full finite-element contact model
              wheel_stage3.py -> wheel_objective.py -> wheel_adjoint.py
              -> wheel_fem.py -> wheel_wheel.py                ~hours
                    |
                    |   best_solution.json (promoted)
                    v
   EXPORT     watertight B-rep solid, filleted, in STEP
              wheel_step_export.py (CadQuery/OpenCascade)      ~minutes
                    |
                    v
              export/wheel.step + wheel_step_manifest.json
\end{lstlisting}

The two search stages exist because they answer different questions at different
prices. The beam model in \code{wheel\_fea.py} scores one design in about 0.25\,ms,
which is what makes a 300-population $\times$ 300-generation genetic search possible.
The finite-element model in \code{wheel\_fem.py} scores one design in roughly 18
seconds per load phase, eight phases per evaluation---but it can see things the beam
model is structurally blind to.

The single most important measured finding in the tree is that \textbf{the beam model
is an accurate spoke model and a blind wheel model}: four GA winners whose losses agree
to 1.2\% and whose beam deflections agree to 0.5\% have real axle drops of 0.71, 1.67,
1.86 and 2.69\,mm---a factor of 3.8 (\at{wheel\_wheel.py}{522--548}). The mechanism is
named there too: the beam model integrates the spoke over its whole hub-to-rim span,
but the built part is welded into its rings over two arcs, so only the middle 67--87\%
of the spoke is a flexure at all.

\subsection{Why the repository looks the way it does}

Three structural oddities will hit you immediately, and all three are deliberate.

\paragraph{Two virtual environments.}
\code{.venv-opt} (numpy/scipy/JAX/pygad) runs the optimizer and the FEA;
\code{.venv-cad} (CadQuery/OpenCascade) runs the STEP exporter. CadQuery pins an
incompatible OCC stack and cannot coexist with the rest (\code{requirements-cad.txt}).
The two interpreters share \textbf{only source files and JSON}. This forces a hard
constraint enforced by \code{tests/test\_import\_hygiene.py}: \code{wheel\_fea.py},
\code{wheel\_geometry.py} and \code{project\_paths.py} must import with numpy
alone---no JAX at module scope---because the CAD environment imports them across the
interpreter boundary (\at{wheel\_step\_export.py}{71--84}). \code{wheel\_geometry.py}
solves this by taking its array module as an argument (\code{xp=np} or
\code{xp=jax.numpy}) rather than importing one.

\paragraph{Extraordinary comment density.}
A docstring in this tree is not documentation, it is \emph{evidence}.
\at{wheel\_objective.py}{165--297} is 130 lines of comment justifying three
floating-point constants, including the measurement that falsified the previous values.
\code{CLAUDE.md} explicitly carves out an exception to its own simplicity rule for
exactly this reason: ``a docstring here carries the measurement that justifies a
constant, and deleting it loses evidence, not verbosity.''

\paragraph{The plan files.}
\code{PLAN.md} is 778\,KB and is the project record, organised into numbered sections
(\S1\,\ldots\,\S92). Eight sibling \code{*\_PLAN.md} files are open work arcs. Comments
across the tree cite these by number. Six closed arc files were deleted on 2026-08-16
and about seventy citations to them remain dangling on purpose---\at{PLAN.md}{53--77}
is the redirect table.

% ===========================================================================
\section{Mathematical foundations}
\label{sec:math}
% ===========================================================================

This section teaches the specific mathematics the code implements, in the order the
data flows. Notation is introduced as it appears.

\subsection{The genome: what a design \emph{is}}
\label{sec:genome}

A design is a point in $\mathbb{R}^{14}$. The gene names, in a fixed order that is a
hard contract (\at{wheel\_genome.py}{29--36}):

\begin{lstlisting}
cx1 cy1 cx2 cy2 cx3 cy3 cx4 cy4     four interior Bezier control points (mm)
t0  t1  t2  t3                       four thickness nodes (mm)
R_hub  R_rim                         two fillet radii (mm)
\end{lstlisting}

Each gene has a box constraint (\at{wheel\_fea.py}{261--310}). Stage 3 works in a
\emph{normalised unit box} $z \in [0,1]^{14}$ rather than in millimetres, because
\code{cy*} spans 64\,mm while \code{R\_rim} spans 2.5\,mm and one gradient step length
cannot mean anything in raw units. The map is affine (\at{wheel\_genome.py}{76--88}):
%
\begin{align}
  z &= \frac{g - \ell}{h - \ell} && \text{(\code{normalize})} \\
  g &= \ell + z\,(h - \ell)      && \text{(\code{denormalize})}
\end{align}
%
with $\ell$ and $h$ the per-gene bounds, and the chain rule
$\partial L/\partial z = (\partial L/\partial g) \odot (h - \ell)$ lives in exactly one
place, \code{wheel\_objective.objective} (\at{wheel\_objective.py}{1376--1379}).

Provenance is enforced structurally. \code{genome\_hash} (\at{wheel\_genome.py}{118})
is
%
\begin{equation*}
  \mathrm{sha256}\bigl(\texttt{json.dumps(sorted(genes.items()))}\bigr)[{:}7],
\end{equation*}
%
so \textbf{adding one key inside \code{genes} changes the hash of every genome ever
recorded}. \code{save\_record} raises if \code{genes} is not exactly the 14 canonical
keys; everything else goes into top-level blocks. The shipped genome's hash is
\code{09e8188}.

\subsection{B\'ezier centerlines}
\label{sec:bezier}

Each spoke's centerline is a \emph{degree-5 B\'ezier curve} in the spoke's own local
frame, running from the hub at $(0,0)$ to the rim at $(S,0)$, with
$S = 48.5 - 12.7 = 35.8$\,mm.

A B\'ezier curve of degree $n$ over control points $P_0,\dots,P_n$ is
%
\begin{equation}
  C(u) \;=\; \sum_{k=0}^{n} B_{n,k}(u)\,P_k,
  \qquad
  B_{n,k}(u) \;=\; \binom{n}{k}\,(1-u)^{\,n-k}\,u^{k},
\end{equation}
%
where the $B_{n,k}$ are the \emph{Bernstein basis polynomials}. Two facts matter here.

\begin{enumerate}
\item \textbf{Linearity in the control points.} For a fixed sample grid
  $u = \mathrm{linspace}(0,1,N)$ the curve is $\mathrm{curve} = B\,P$ with $B$ a
  constant $(N,6)$ matrix. That is why the gradient path from genes to node positions
  is cheap---the map is a matrix multiply. The basis is \code{lru\_cache}d because it
  depends only on $N$, which removed a measured 21\% of a fitness evaluation
  (\at{wheel\_geometry.py}{83--97}).

\item \textbf{Endpoint tangency.} The derivative curve (the \emph{hodograph}) of a
  degree-$n$ B\'ezier is a degree-$(n-1)$ B\'ezier over the scaled forward differences
  $n\,\Delta P_k$. At $u = 0$ this gives $C'(0) = n\,(P_1 - P_0)$, so \textbf{the
  endpoint tangent is exactly the first control-polygon edge}. Since $P_0$ is locked at
  the origin, the hub tangent direction is exactly $(\code{cx1}, \code{cy1})$.
  \code{wheel\_geometry.forward\_difference\_matrix} (\code{:117}) builds the difference
  operator and \code{bezier\_tangent} (\code{:148}) applies it. This closed form is what
  makes the \code{arrival} barrier free to evaluate (\S\ref{sec:arrival}) and is what
  \code{studies/study\_arrival\_cap.py} exploits to sweep one angle while moving no
  other gene.
\end{enumerate}

\code{bezier\_curvature} (\at{wheel\_geometry.py}{164}) gives the signed curvature
analytically,
%
\begin{equation}
  \kappa(u) \;=\; \frac{x'y'' - y'x''}{\bigl(x'^2 + y'^2\bigr)^{3/2}},
  \label{eq:curvature}
\end{equation}
%
from the first and second hodographs. This is needed for the fold constraint below.

\subsection{The offset band: turning a curve into a part}
\label{sec:band}

The spoke is the centerline \emph{thickened along its normal by a variable amount}.
Define normalised arc length $s \in [0,1]$ and a transverse coordinate
$\eta \in [-1,+1]$. Then
%
\begin{equation}
  x(s,\eta) \;=\; C(s) \;+\; \eta\,\frac{t(s)}{2}\,\hat{N}(s),
  \label{eq:offsetband}
\end{equation}
%
where $\hat N$ is the unit normal (the tangent rotated $+90^\circ$) and $t(s)$ is the
local thickness. This one formula is \code{wheel\_geometry.offset\_band}
(\code{:290}), and it does double duty:

\begin{itemize}
\item at \code{n\_across = 1} it returns the two flanks---the outline the STEP exporter
  builds from;
\item at \code{n\_across = 6} it returns a structured quadrilateral mesh block.
\end{itemize}

\textbf{The meshed part and the exported part are therefore the same surface by
construction}, not by two pipelines agreeing (\at{wheel\_mesh.py}{14--21}).

\paragraph{Thickness} is piecewise-linear over three equal arc-length zones with nodes
$t_0 \ldots t_3$ at $s = 0, \tfrac13, \tfrac23, 1$. The original numpy implementation
used boolean masks; JAX has no equivalent, because \code{s[mask]} has a data-dependent
shape. \code{thickness\_at\_arc\_length} (\at{wheel\_geometry.py}{217}) rewrites it
branch-free as a base value plus three clipped ramps:
%
\begin{equation}
  t(s) \;=\; t_0 \;+\; \sum_{i=0}^{2} \bigl(t_{i+1} - t_i\bigr)\,
             \operatorname{clip}\!\left(\frac{s - b_i}{b_{i+1} - b_i},\,0,\,1\right).
  \label{eq:taper}
\end{equation}
%
The equivalence argument is in the docstring: inside zone $i$, ramps $j < i$ are
saturated at 1 and telescope to $t_i - t_0$, ramp $i$ contributes the interpolation
fraction, and ramps $j > i$ clip to 0. This is exactly \code{np.interp}, and the
docstring records that dropping the original's spurious $+10^{-12}$ made the function
hit its own interpolation nodes exactly.

\subsection{The fold constraint: where an offset curve destroys itself}
\label{sec:fold}

This is the first piece of genuinely non-obvious geometry.

If you offset a curve outward by a distance $d$, and the curve's local \emph{radius of
curvature} $\rho = 1/|\kappa|$ drops below $d$, the offset passes \emph{through the
centre of curvature} and turns inside out. The flank crosses itself; the region has
negative area; the corresponding mesh element is inverted; the STEP wire
self-intersects. The part described by such a genome \textbf{does not exist}, and yet
the Castigliano beam model will happily return a deflection for it.

The signed clearance is (\code{wheel\_geometry.self\_intersection\_margin},
\code{:421}):
%
\begin{equation}
  \mathrm{margin} \;=\; \min_{s}\;\left(\frac{1}{|\kappa(s)|} \;-\; \frac{t(s)}{2}\right).
  \label{eq:foldmargin}
\end{equation}
%
Negative means folded. This is exposed as a \emph{smooth signed scalar} precisely so it
can serve as a barrier term rather than be discovered as a crash.

The threshold is not zero, and the reason is measured
(\at{wheel\_geometry.py}{355--374}). Over 2001 feasible genomes sampled by
\code{studies/study\_mesh\_quality.py}:

\begin{table}[h]
\centering
\small
\begin{tabular}{@{}rrrr@{}}
\toprule
$\mathrm{margin} >$ & missed bad meshes & false alarms & designs kept \\
\midrule
0.00 & 6 & 87 & 1187 \\
0.05 & 2 & 116 & 1154 \\
\textbf{0.10} & \textbf{0} & 140 & 1128 \\
0.20 & 0 & 196 & 1072 \\
\bottomrule
\end{tabular}
\caption{Why \code{MIN\_FOLD\_MARGIN\_MM} is 0.1 and not 0.}
\end{table}

\code{MIN\_FOLD\_MARGIN\_MM = 0.1} is the smallest threshold with zero misses, and it
keeps 95\% of the designs a zero threshold would. Roughly \textbf{40\% of the genomes
the other constraints admit} have a folded region (\at{wheel\_fea.py}{653--671})---this
is not an exotic failure.

\subsection{The arrival angle, and why its sense is backwards}
\label{sec:arrival}

The centerline endpoints are \emph{locked} on the hub and rim circles, and the end
cross-section is normal to the centerline. Define the \emph{arrival angle} as the angle
between the centerline and the ring's \emph{tangent} at that endpoint: $0^\circ$
tangential, $90^\circ$ radial. Closed form
(\code{wheel\_wheel.arrival\_angles}, \code{:380}), using the hodograph endpoint fact
from \S\ref{sec:bezier} and the fact that the endpoint's own unit vector \emph{is} the
local radial direction:
%
\begin{equation}
  \mathrm{arrival} \;=\; \arcsin\bigl|\,\hat d \cdot \hat r\,\bigr|.
\end{equation}

The intuition is that a near-\emph{tangent} arrival is the hard case---and for the CAD
fillets it is. For the \emph{mesh} it is exactly backwards. Because the end
cross-section is normal to the centerline, a near-tangent arrival makes the
cross-section nearly \emph{radial}, so it meets the ring arc at about $80^\circ$ and
the junction block is beautifully shaped. A near-\emph{radial} arrival lays the
cross-section \emph{along} the arc and the corner collapses. Measured over 58 feasible
genomes, the correlation between arrival angle and junction minimum scaled Jacobian is
$-0.93$ at the hub and $-0.91$ at the rim:

\begin{table}[h]
\centering
\small
\begin{tabular}{@{}lrrrr@{}}
\toprule
arrival (deg) & 4.9 & 7.2 & 64.8 & 82.2 \\
min.\ scaled Jacobian & 0.806 & 0.804 & 0.415 & 0.043 \\
\bottomrule
\end{tabular}
\end{table}

Past about $80^\circ$ both flanks lie \emph{outside} the ring circle and the spoke
touches its ring at a single point. That is a \textbf{hinge, not a weld}---and since
the beam model assumes a moment connection at both ends
(\code{BOUNDARY\_CONDITION = "fixed\_guided"}), it would be solving a structure the part
does not have. The threshold sweep over 200 genomes is sharp: the worst genome failing
$\mathrm{minSJ} > 0.2$ arrives at $70.6^\circ$, and every threshold at or below
$70^\circ$ gives zero misses. \code{MAX\_ARRIVAL\_DEG = 65.0} takes $5.6^\circ$ of
margin (\at{wheel\_wheel.py}{355--377}).

\subsection{Barriers: how constraints become part of a smooth objective}
\label{sec:barriers}

Every constraint in this tree is a \emph{quadratic soft barrier}
(\at{wheel\_fea.py}{579}; the \code{jnp} version at \at{wheel\_objective.py}{414}):
%
\begin{equation}
  \mathrm{soft\_barrier}(v, w) \;=\; w \cdot \max(0, v)^2 .
  \label{eq:softbarrier}
\end{equation}
%
Zero when satisfied, rising quadratically when violated. Quadratic rather than linear
so that it is $C^1$ at the knee---which is exactly what a finite-difference plateau
needs, and what a gradient-based line search needs in order to avoid a kink. This
replaced a hard $-500$ cliff that a gradient-free search could not descend.

The repository is careful about one distinction that is \emph{not} derivable from the
weight table (\at{wheel\_objective.py}{394--407}):

\begin{itemize}
\item a \textbf{barrier term} answers ``may this design ship?''---its only admissible
  value is 0, and it never trades against mass;
\item an \textbf{objective term} answers ``how good is it?''
\end{itemize}

Conflating them caused a real defect: a run reported its lowest-\emph{loss} iterate,
which was one that violated the hub fillet cap slightly and bought a lot of mass for
it. Fifty-five of its 101 iterates were feasible and it promoted an infeasible one. The
fix is \code{wheel\_stage3.selection\_key} (\code{:195}), a three-tier lexicographic
sort described in \S\ref{sec:selectionkey}.

\subsection{The beam model: the force method (Castigliano)}
\label{sec:castigliano}

Stage 1/2's physics is \code{wheel\_fea.generalized\_spoke\_mechanics} (\code{:398}).
It solves one curved spoke as a \emph{statically indeterminate} beam by the force
method.

Set the local frame with the hub root at $(0,0)$ and the rim tip at $(S,0)$. The radial
service load $F$ acts at the tip along $-x$. If the tip were free, the internal moment
and axial force at a station would be determined by statics alone. It is not free: the
spoke is fused into a stiff rim ring, which restrains tip rotation and tangential
motion. So two \emph{redundants} are carried at the tip:

\begin{itemize}
\item $M_0$ --- an end moment, restraining tip rotation;
\item $H$ --- a tangential force, restraining tip motion along $+y$.
\end{itemize}

Cutting a free body from station $s$ to the tip gives the internal actions
%
\begin{align}
  M(s) &= M_0 + H\,(S - x) - F\,y, \\
  N(s) &= -F\cos\theta + H\sin\theta,
  \qquad (\cos\theta, \sin\theta) = \left(\tfrac{dx}{ds}, \tfrac{dy}{ds}\right),
\end{align}
%
with $N$ \textbf{negative in compression} (v2.0 had this backwards, which is why its
buckling check never fired). Write these as linear combinations of unit-load fields:
%
\begin{align}
  M &= M_0\,m_{M_0} + H\,m_H + F\,m_F, &
      m_{M_0} &= 1, & m_H &= S - x, & m_F &= -y,\\
  N &= M_0\,n_{M_0} + H\,n_H + F\,n_F, &
      n_{M_0} &= 0, & n_H &= \sin\theta, & n_F &= -\cos\theta.
\end{align}

The \emph{complementary strain energy} of a beam carrying moment and axial force is
%
\begin{equation}
  U \;=\; \int \frac{M^2}{2EI}\,ds \;+\; \int \frac{N^2}{2EA}\,ds .
\end{equation}
%
\textbf{Castigliano's second theorem} says $\partial U/\partial X = 0$ for any redundant
$X$ whose corresponding displacement is restrained to zero. Applying it to $M_0$ and
$H$ gives a $2\times2$ linear system in the \emph{flexibility coefficients}
%
\begin{equation}
  a_{ij} = \int \frac{m_i m_j}{EI}\,ds + \int \frac{n_i n_j}{EA}\,ds,
  \qquad
  b_{i}  = \int \frac{m_i m_F}{EI}\,ds + \int \frac{n_i n_F}{EA}\,ds,
\end{equation}
%
\begin{equation}
  \begin{bmatrix} a_{11} & a_{12} \\ a_{12} & a_{22} \end{bmatrix}
  \begin{bmatrix} M_0 \\ H \end{bmatrix}
  \;=\; -F \begin{bmatrix} b_1 \\ b_2 \end{bmatrix},
  \label{eq:flexibility}
\end{equation}
%
which the code assembles as segment sums over the 600-point discretization
(\at{wheel\_fea.py}{495--520}). A tiny relative ridge $+10^{-12}(1 + |a_{ii}|)$ is added
to the diagonal, because a near-straight, very thin spoke makes the system
ill-conditioned and a genetic algorithm will propose exactly that.

The tip deflection along the load direction is then Castigliano's \emph{first}
application, $\partial U/\partial F$:
%
\begin{equation}
  \delta \;=\; \left|\; \sum \frac{M\,m_F}{EI}\,\Delta s
                 \;+\; \sum \frac{N\,n_F}{EA}\,\Delta s \;\right|.
\end{equation}
%
The comment at \at{wheel\_fea.py}{522--525} makes an argument worth understanding: by
the \emph{envelope theorem}, the redundants sit at stationary values of $U$, so
differentiating with respect to $F$ while \emph{holding them fixed} is exact---the
chain terms through $dM_0/dF$ and $dH/dF$ vanish. No extra work is needed.

The legacy \code{bc="cantilever"} sets $M_0 = H = 0$, which collapses algebraically to
$M = -Fy$ and $N = -F\cos\theta$---term for term the v2.0 integrands. That equivalence
is the regression test. The physical consequence of the change: for a straight beam,
cantilever gives $FL^3/3EI$ and fixed-guided gives $FL^3/12EI$, so \textbf{v2.0
over-predicted compliance by a factor of four}.

Three corrections ride on top.

\begin{itemize}
\item \textbf{Peterson stress concentration} (\at{wheel\_fea.py}{360}). At a fillet in a
  stepped beam in bending,
  $K_t \approx 1 + C\,(t/2R)^{0.65}$, clamped to $[1.0, 3.5]$. The peak bending stress
  at each end is multiplied by its local $K_t$. This same formula reappears---%
  differentiated rather than frozen---as the FEA-era stress constraint
  (\S\ref{sec:stress}).

\item \textbf{Winkler--Bach curved-beam correction} (\at{wheel\_fea.py}{376}). On a
  curved member the neutral axis shifts toward the centre of curvature and inner-fibre
  stress is higher than straight-beam theory says. Simplified: $k_b = 1 + h/(4R)$,
  clipped to $[1.0, 2.5]$, significant only when $R < 4h$.

\item \textbf{Euler buckling as an equivalent column} (\at{wheel\_fea.py}{543--556}):
  %
  \begin{equation}
    P_e \;=\; \frac{\pi^2\,EI_{\text{eff}}}{(K_e L)^2},
    \qquad
    EI_{\text{eff}} \;=\; \frac{L}{\sum \Delta s / EI},
  \end{equation}
  %
  where $EI_{\text{eff}}$ is the \emph{compliance-weighted (harmonic) mean} rigidity---%
  the standard equivalent-column treatment for a tapered member. v2.0 applied this per
  discretization segment, where $L \approx 0.06$\,mm made $1/L^2$ astronomical and the
  check could never fire. The docstring is honest that for an in-plane-curved member
  Euler is a proxy: the genuinely competing modes are snap-through and
  lateral-torsional buckling.
\end{itemize}

\subsection{The genetic algorithm}

\code{wheel\_fea.evaluate\_design} (\code{:699}) sums nine weighted terms into a scalar
loss; \code{pygad\_fitness} returns its negation because PyGAD maximises. The GA itself
(\at{wheel\_fea.py}{1310}\,ff.) is population 300 (at least $20\times$ the gene count),
tournament selection with $k=5$, uniform crossover, elitism 5, and an \emph{adaptive
Gaussian mutation} whose $\sigma$ starts at 0.25 of the gene range and halves halfway
through the run (\code{adaptive\_gaussian\_mutation}, \code{:881}).

Two of the loss terms are worth naming, because they are the \emph{mesh's} constraints
imported into a search that predates the mesh: \code{fold} (\S\ref{sec:fold}) and
\code{arrival} (\S\ref{sec:arrival}), both closed-form and therefore free next to the
Castigliano solve, and both fixing cases where \textbf{the model itself is invalid}
rather than merely where the design is poor.

\subsection{Finite elements: write the energy, derive everything else}
\label{sec:fem}

\code{wheel\_fem.py} is the FEA kernel, and its organizing principle is stated in the
first paragraph of its docstring: \textbf{\code{element\_energy} is the only physics in
the file.} Internal force and tangent stiffness are its gradient and Hessian, obtained
by \code{jax.grad} and \code{jax.hessian} and \code{vmap}ped over elements
(\at{wheel\_fem.py}{258--278}):
%
\begin{equation}
  f_e \;=\; \nabla_{u}\,\Pi_e,
  \qquad
  K_e \;=\; \nabla^2_{u}\,\Pi_e .
\end{equation}
%
There is no B-matrix, no hand-differentiated stress recovery, and no separately coded
tangent that can drift out of step with the residual. That entire class of bug is
absent by construction, and the rigid-body test becomes exact rather than a
cancellation: if the energy is frame-indifferent, so is everything derived from it.

\paragraph{Elements.}
Isoparametric quadrilaterals: Q4 (bilinear, 4 nodes) or Q9 (biquadratic, 9 nodes). Q9
is the production choice because \textbf{a Q4 mesh shear-locks in bending}---it reports
the structure as stiffer and safer than it is, which is catastrophic for a 2\,mm
flexure. Full Gauss integration ($2\times2$ for Q4, $3\times3$ for Q9); reduced
integration is the usual dodge for Q4 locking but buys hourglass modes
(\at{wheel\_fem.py}{144--160}).

\paragraph{The Jacobian, and one index that matters.}
In \code{element\_energy} (\code{:235}):

\begin{lstlisting}[style=py]
J      = jnp.einsum("gnk,ni->gki", dN, Xe)      # J[g,k,i] = dx_i/dxi_k
detJ   = J[:,0,0]*J[:,1,1] - J[:,0,1]*J[:,1,0]
dNdx   = jnp.einsum("gnk,gik->gni", dN, jnp.linalg.inv(J))
grad_u = jnp.einsum("ni,gnj->gij", ue, dNdx)
\end{lstlisting}

\code{J} here is the \emph{transpose} of the usual Jacobian matrix, so
$\partial \xi_k/\partial x_i$ is \code{inv(J)[i,k]}, not \code{inv(J)[k,i]}. Getting it
backwards is invisible on an axis-aligned rectangle ($J$ is diagonal) and silently
wrong on every curved element---which is exactly the geometry this project meshes. The
distorted-mesh patch test is what catches it.

\paragraph{Constitutive law.}
St.\ Venant--Kirchhoff strain energy density
(\code{wheel\_fem.\_energy\_density}, \code{:221}):
%
\begin{equation}
  W(\varepsilon) \;=\; \frac{\lambda}{2}\,\bigl(\operatorname{tr}\varepsilon\bigr)^2
                  \;+\; \mu\,\bigl(\varepsilon : \varepsilon\bigr),
  \label{eq:svk}
\end{equation}
%
with Lam\'e constants from \code{lame()} (\code{:190}). \emph{Plane stress} uses the
reduced $\lambda = E\nu/(1-\nu^2)$; plane strain uses
$\lambda = E\nu/\bigl((1+\nu)(1-2\nu)\bigr)$. The choice is deliberate and recorded in
the result dictionary: plane stress makes a uniaxial state give $\sigma = E\varepsilon$
exactly, so beam bending gives $EI = E w t^3/12$ with \emph{plain} $E$---apples to
apples with the Castigliano model, which is what makes the beam-versus-FE comparison a
check rather than an argument. The real 22.4\,mm-wide part is closer to plane
\emph{strain} (effective modulus $E/(1-\nu^2) = 1.14E$), and that 14\% is larger than
most effects the project chases, so \code{plane="strain"} exists and is never picked
silently (\at{wheel\_fem.py}{49--61}).

\subsection{Two kinematics, five lines apart}
\label{sec:kinematics}

\code{\_strain} (\at{wheel\_fem.py}{207}) returns either
%
\begin{align}
  \text{linear:}\quad \varepsilon &= \tfrac12\bigl(\nabla u + \nabla u^{\mathsf T}\bigr),
  \\
  \text{svk:}\quad E &= \tfrac12\bigl(F^{\mathsf T}F - I\bigr),
  \qquad F = \nabla u + I .
\end{align}

\textbf{Green--Lagrange is exactly frame-indifferent.} Under a rigid rotation
$u = (R - I)x$ we have $F = R$, so $E = \tfrac12(R^{\mathsf T}R - I) = 0$ identically:
a rigid rotation of \emph{any} magnitude stores zero energy. The linear strain does not
have this property, and that is the entire content of the finite-rotation test. Because
the two share the same \code{\_energy\_density}, they cannot disagree about the
material.

There is a trap the file names explicitly (\at{wheel\_fem.py}{980--1010}): at $u = 0$
the SVK Hessian \emph{equals} the linear one exactly. So routing an \code{svk} problem
through \code{solve\_linear} returns a linear answer silently. The fix is
\code{solve()} (\code{:1260}), which dispatches on \code{prob.nonlinear}.

How much does it matter? \code{studies/study\_gnl.py} measures it: at service load the
wheel is \textbf{23.16\%} softer under SVK than under linear kinematics for the shipped
genome---and, more decisively, that correction is \emph{not} a property of the wheel
that can be applied once. Held at a fixed axle drop by scaling the load per genome, it
still ranges over roughly an order of magnitude across feasible designs. So the
nonlinear solve has to be in the loop. \code{wheel\_stage3.py --kinematics} defaults to
\code{svk} since \code{PLAN} \S32; \code{wheel\_fem}'s own kernel defaults stay
\code{linear} on purpose so M4's committed reports stay reproducible.

\subsection{Contact: the ground is not a pressure you assume}
\label{sec:contact}

Early full-wheel work applied the ground load as an assumed elliptical pressure over an
assumed $3^\circ$ patch (\code{\_ground\_traction}, \at{wheel\_fem.py}{1462}). Two
choices there are worth noting: the traction is \textbf{vertical, not radial} (for
frictionless contact against a rigid \emph{flat} ground the bodies share a tangent
plane, so the traction is along the ground's normal---using the rim's radial normal
produces a spurious horizontal resultant), and it is \textbf{elliptical, not uniform}
(a uniform patch has a pressure step at its edges, and a step puts a stress singularity
into a model whose peak stress is a reported number).

But the patch half-angle was a free parameter, and sweeping it moved the axle drop by
9.7\% over a $12\times$ range---larger than the entire geometric-nonlinearity
correction. So M6 replaced it with \emph{penalty contact} against a rigid frictionless
ground.

\code{RigidGroundContact} (\at{wheel\_fem.py}{652}) adds a potential to the total
energy. For a boundary segment with gap $g = (x_y + u_y) - y_{\text{ground}}$, the
penalty potential density is $\varepsilon_n\,\varphi(-g)$ where $\varphi$ is a smoothed
positive part (\code{\_penalty\_phi}, \code{:592}). The sharp branch is
$\varphi(z) = \tfrac12\langle z\rangle^2$, which is $C^1$; the smoothed branch is a
cubic matching value, slope \emph{and} curvature at both $z = 0$ and $z = d$, so
$\varphi''$ is continuous everywhere:
%
\begin{equation}
  \varphi(z) \;=\;
  \begin{cases}
    0, & z \le 0,\\[2pt]
    \dfrac{z^3}{6d}, & 0 < z < d,\\[6pt]
    \dfrac{z^2}{2} - \dfrac{dz}{2} + \dfrac{d^2}{6}, & z \ge d .
  \end{cases}
  \label{eq:penalty}
\end{equation}
%
The same \code{grad}/\code{hessian} machinery produces the contact force and contact
stiffness, so they cannot disagree with the potential.

\paragraph{The problem is displacement-driven, and that is not a convenience.}
With a rigid flat ground indenting by $\delta$, the rise of the wheel's lowest point
\emph{is} $\delta$, so the axle drop is an \textbf{input} and the wheel load is the
\textbf{output} (\at{wheel\_fem.py}{1735--1751}). A force-driven contact problem is
singular at its own starting point---before contact establishes there is no pressure at
all. \code{solve\_wheel\_contact} (\code{:1867}) then inverts this with a \emph{secant
iteration} to answer the design question (``what is the drop at 66.72\,N?''), each
solve warm-starting from the previous.

The measured payoff: the real patch is a half-angle of about \textbf{$0.47^\circ$}, not
$3.0^\circ$---the assumed patch was six times too wide---and yet the axle drop moves
only 1.3\%, because the drop is dominated by spoke and rim bending rather than by how
the last few newtons are spread (\at{studies/study\_contact.py}{14--30}).

\subsection{Constraints by master--slave elimination}
\label{sec:dofmap}

\code{DofMap} (\at{wheel\_fem.py}{828}) expresses every boundary condition as one
sparse matrix,
%
\begin{equation}
  u_{\text{full}} \;=\; T\,u_{\text{reduced}} \;+\; u_{\text{prescribed}},
\end{equation}
%
and the reduced system $T^{\mathsf T} K T$ is symmetric positive definite whenever the
full one is. This single mechanism handles homogeneous and inhomogeneous Dirichlet
conditions, \emph{skew} constraints (\code{constrain\_direction}: enforce
$u \cdot \hat d = 0$, leaving the perpendicular free---this is how ``plane sections
remain plane but may stretch laterally'' is expressed without rotating the system into
a local frame), and \textbf{rigid ties}:
%
\begin{align}
  u_x &= U_x - \Theta\,(y - y_c), \\
  u_y &= U_y + \Theta\,(x - x_c),
\end{align}
%
which tie a whole cross-section to a 3-DOF rigid body. That is how the spoke's tip
boundary condition is made to mean exactly what beam theory assumes: plane sections
remain plane \emph{and normal}. The two redundants of \S\ref{sec:castigliano} reappear
here as reactions---$M_0$ is the reaction of $\Theta = 0$, and $H$ the reaction of the
transverse constraint.

\code{finalize()} asserts that \textbf{every DOF is claimed exactly once}: an unclaimed
DOF is a free floating node (singular matrix), and a doubly-claimed one silently drops
the first constraint.

\subsection{Newton with an energy line search}
\label{sec:newton}

\code{solve\_nonlinear} (\at{wheel\_fem.py}{1116}) is a Newton--Raphson loop with load
continuation and Armijo backtracking. One detail is genuinely instructive:

\begin{lstlisting}[style=py]
slope = float(r @ du)
if slope >= 0.0:
    raise NewtonDivergedError("...the tangent is not positive definite,
                               i.e. this equilibrium is unstable")
\end{lstlisting}

Because \code{r} really \emph{is} $\nabla\Pi$ rather than an approximation of it,
\code{slope} is exact. So a non-negative value is not line-search noise---it says the
tangent lost positive definiteness, i.e.\ \textbf{this equilibrium is unstable}. Saying
that is far more useful than letting the backtrack loop grind $\alpha$ down to nothing.
This is also, as \at{wheel\_stage3.py}{76--85} notes, the only buckling signal the
repository owns until the Euler proxy is replaced.

Convergence fires on \emph{either} the relative reduced residual
$\lVert r \rVert / \lVert f_{\text{ref}} \rVert \le \mathrm{tol}$ \emph{or} the energy
increment $|du \cdot r|$ falling to $\mathrm{tol}_{\text{energy}}$ of the first step's
--- the second because at low load it is the only criterion attainable.

\subsection{Stress: why the maximum is not a number}
\label{sec:stress}

The peak von Mises stress in this wheel \textbf{diverges under mesh refinement}:
$31.02 \to 41.54 \to 48.47$\,MPa over coarse/medium/fine. That is not a solver defect.
The spoke/ring junction is a \emph{re-entrant corner}, and the elastic field at a
traction-free wedge of material angle $\omega$ behaves as $\sigma \sim r^{\lambda - 1}$
where $\lambda$ is the smallest root in $(0,1)$ of the \emph{Williams eigenvalue
equation} (mode I):
%
\begin{equation}
  \sin(\lambda\omega) \;+\; \lambda\,\sin(\omega) \;=\; 0 .
  \label{eq:williams}
\end{equation}
%
\at{studies/study\_corner\_singularity.py}{132} solves this by bisection and verifies it
against the two cases with known answers: $\omega = 360^\circ$ (a crack) gives exactly
$0.5$, and $\omega = 270^\circ$ gives the textbook $0.5445$. The measured junction
wedges here are $295^\circ$--$356^\circ$, all re-entrant, so $\lambda - 1 < 0$ and the
stress is unbounded at the corner. \textbf{The maximum is a property of the mesh, not
of the wheel.}

This has three consequences the code carries.

\paragraph{(a) The optimizer cannot use \texttt{max}.}
Its derivative is a delta on whichever Gauss point happens to be highest, so it jumps
discontinuously the moment the argmax moves---which, on a rolling wheel with a contact
patch sweeping the rim, is every few steps. The smooth surrogate is a
\emph{volume-weighted $p$-norm} (\code{wheel\_adjoint.\_qoi\_pnorm\_stress},
\code{:313}):
%
\begin{equation}
  \sigma_{pn} \;=\;
  \left(\frac{\sum_g V_g\,\mathrm{vm}_g^{\,p}}{\sum_g V_g}\right)^{\!1/p},
  \qquad
  V_g \;=\; w_g\,\det J_g\,\cdot\,\mathrm{width}.
  \label{eq:pnorm}
\end{equation}
%
Volume weighting is essential: an unweighted $p$-norm scores the \emph{mesh}, because
refining the rim adds Gauss points in a lightly stressed region and drags the average
down---the optimizer could reduce the objective by changing \code{cfg}. With volume
weights the sum is a quadrature of an integral and is mesh-convergent.

\paragraph{(b) Rescaling a $p$-norm to the maximum cannot work.}
The constraint used to be $c\,\sigma_{pn} \le \mathrm{allowable}$ with
$c = \max/\sigma_{pn}$ measured and frozen within a step. The freeze was mathematically
sound---the $p$-norm is positively homogeneous of degree 1, so
$d(c\,\sigma_{pn}) = c\,d(\sigma_{pn})$ \emph{if and only if} $c$ did not move---but the
quantity was not. \code{make m8bi6} swept ten exponents off one set of solves and found
the two factors converge in \textbf{disjoint ranges of $p$}: the $p$-norm for
$p \le 6$, the ratio $c$ only for $p \ge 24$, and their product at \emph{no} exponent at
either design. Here $c$ is anchored to a divergent quantity, and a frozen divergent
number is still a divergent number.

\paragraph{(c) So the peak is modelled, not measured.}
The shipped constraint (\at{wheel\_objective.py}{1160--1200}) is
%
\begin{equation}
  K_t(R, t)\;\cdot\;\sigma_{\text{nominal}}(p{=}4)
  \;\;\le\;\; \sigma_{\text{allow}} \;=\; 25.0\ \mathrm{MPa},
  \label{eq:stressconstraint}
\end{equation}
%
where $\sigma_{\text{allow}}$ is \code{ALLOWABLE\_STRESS\_MPA}. It is
one \code{soft\_barrier} per junction, summed---the hub priced on
$(\code{R\_hub}, t_0)$, the rim on $(\code{R\_rim}, t_3)$. Here
$\sigma_{\text{nominal}}$ at $p = 4$ is mesh-convergent (GCI 0.45\% / 0.20\%, observed
order 1.76 / 2.18, comparable to the axle drop's 2.44 / 0.14\%), and
$K_t = 1 + C(t/2R)^{0.65}$ is the same Peterson factor from \S\ref{sec:castigliano},
now \textbf{differentiated by \code{jax.grad} rather than frozen}. Two junctions get two
barriers rather than one aggregated $K_t$, because a hard \code{max} over the two would
zero the gradient of whichever is not currently worst and reintroduce exactly the
argmax kink the $p$-norm exists to blur.

Note the consequence: \code{STRESS\_NOMINAL\_P = 4.0} (\at{wheel\_objective.py}{147}) is
the constraint's exponent, while \code{wheel\_adjoint.STRESS\_PNORM\_P = 30.0}
(\code{:283}) stays the adjoint's documented default so historical records keep their
meaning. Two exponents, two meanings, both pinned by tests.

\subsection{Mesh convergence: Richardson and the GCI}
\label{sec:gci}

Because several key quantities are only meaningful in the limit, the tree does formal
mesh-convergence analysis. Given three solutions $\phi_1$ (coarse) to $\phi_3$ (fine)
at cell sizes $h_1 > h_2 > h_3$, Roache's \emph{observed order of convergence} solves
%
\begin{equation}
  p \;=\; \frac{\bigl|\,\ln|e_{32}/e_{21}| + q(p)\,\bigr|}{\ln r_{21}},
  \qquad
  q(p) \;=\; \ln\!\left(\frac{r_{21}^{\,p} - s}{r_{32}^{\,p} - s}\right),
  \label{eq:roache}
\end{equation}
%
by fixed-point iteration from $q = 0$ (\at{studies/study\_deflection\_gci.py}{128}). The
\emph{Richardson extrapolation} to $h \to 0$ and the \emph{Grid Convergence Index} (a
conservative uncertainty band on the extrapolated value) follow at \code{:166}. The
docstring records a real bug: the first draft used the coarse-pair ratio as the
denominator, because Roache indexes 1 as the \emph{finest} grid while \code{phi} here
runs coarse to fine.

A caution the tree learned the hard way: \textbf{look at the increment ratios before
quoting a ladder.} \code{studies/study\_knee\_rungs.py} shows a five-rung ladder whose
ratio \emph{climbs toward 1} instead of decaying, so every extrapolation lands above
the knee while every measurement lands below it. A rising ratio means the top rung is a
lower bound, not an estimate.

\subsection{The adjoint method}
\label{sec:adjoint}

This is the mathematical centrepiece of Stage 3, and \code{wheel\_adjoint.py}'s module
docstring is the derivation. The setup is a converged contact equilibrium with
%
\begin{align}
  \Pi(c, u_f, y) &= U_{\text{bulk}}(c, u_f) + \Pi_{\text{contact}}(c, u_f, y), \\
  u_f            &= T u_r + u_{\text{pre}}
                    &&\text{($T$ static: topology frozen)}, \\
  r(c, u_r)      &= T^{\mathsf T}\nabla_{u_f}\Pi \;=\; \nabla_{u_r}\Pi \;=\; 0
                    &&\text{at the solution}, \\
  K_r            &= \nabla^2_{u_r}\Pi
                    &&\text{what Newton already assembled},
\end{align}
%
where $c$ is the node coordinates and $y$ the ground position. For a quantity of
interest $Q(c, u_f, y)$, differentiate the equilibrium condition $r = 0$ implicitly.
The result is
%
\begin{align}
  K_r\,\mathrm{adj}_r &= -T^{\mathsf T}\,\frac{\partial Q}{\partial u_f}
     &&\text{one sparse solve; $K_r$ is symmetric},
     \label{eq:adjsolve}\\
  \mathrm{adj}_f &= T\,\mathrm{adj}_r, \\
  \frac{dQ}{dc} &= \left.\frac{\partial Q}{\partial c}\right|_{u}
     + \nabla_c\Bigl[\,\mathrm{adj}_f \cdot \nabla_{u_f}\Pi(c, u_f, y)\,\Bigr],
     \label{eq:adjcore}\\
  \frac{dQ}{d\,\mathrm{genes}} &=
     \mathrm{vjp}_{\code{mesh\_coords}}\!\left(\frac{dQ}{dc}\right).
\end{align}

\textbf{Equation~\eqref{eq:adjcore} is the whole adjoint: one \code{jax.grad} of one
scalar.} There is no separately coded sensitivity operator, for the same reason there
is no separately coded contact stiffness. The cost is one sparse back-solve per
quantity, sharing one factorisation (\code{adjoint\_grads},
\at{wheel\_adjoint.py}{516}, uses \code{scipy.sparse.linalg.factorized} and reuses the
\code{jax.vjp} closure across quantities).

Two things \emph{fall out} rather than being coded:

\begin{itemize}
\item The total contact force is \textbf{exactly} $+\partial\Pi/\partial y_{\text{ground}}$
  --- the penalty potential's derivative with respect to the ground's position \emph{is}
  the resultant it presses with. So \code{\_qoi\_contact\_force} is not a
  reimplementation of the pressure quadrature, and the two agreeing is a free
  correctness check.
\item $d(\text{force})/d(\text{indentation})$ is the same adjoint with the same
  multiplier, differentiated with respect to $y_{\text{ground}}$ instead of $c$.
\end{itemize}

\paragraph{Load control, done as an implicit function.}
Stage 3 loads to a \emph{force}, not an indentation, so every quantity is evaluated at
$\delta^*(p)$ solving $F(\delta^*, p) = F_{\text{service}}$. The naive move---%
differentiate the secant iteration---would put the secant's \emph{stopping tolerance}
into the answer, which is a property of the stopping rule and not of the wheel. Instead
the quotient rule for implicit functions is used directly
(\code{service\_qoi\_value\_and\_grad}, \at{wheel\_adjoint.py}{645}):
%
\begin{align}
  \frac{d\delta^*}{dp}
    &= -\,\frac{\bigl(\partial F/\partial p\bigr)\big|_{\delta}}
               {\bigl(\partial F/\partial \delta\bigr)\big|_{p}},
  \label{eq:implicit}\\[4pt]
  \left.\frac{dQ}{dp}\right|_{\text{total}}
    &= \left.\frac{\partial Q}{\partial p}\right|_{\delta}
       \;+\; \frac{\partial Q}{\partial \delta}\cdot\frac{d\delta^*}{dp}.
  \label{eq:coupling}
\end{align}
%
Numerator and denominator both come from the single \code{contact\_force} adjoint at
the converged state. The coupling term is \textbf{not a correction}: measured at the
shipped genome, dropping it does not lengthen the stress gradient, it \emph{reverses}
it ($+11.22$ against a true $-2.59$). Keeping only the frozen term leaves a gradient
with the right direction and the wrong length---and no line search reports that as an
error; it reports it as a hard problem.

There is also a guard worth quoting, because it encodes physics as a sanity check:

\begin{lstlisting}[style=py]
if not np.isfinite(dF_ddelta) or dF_ddelta <= 0.0:
    raise RuntimeError("a wheel that gets softer the harder it is pressed
                        is not a load case, it is a solver failure")
\end{lstlisting}

\subsection{Structured multi-block meshing}
\label{sec:meshing}

The mesh is not produced by an external mesher, and that is a load-bearing decision:
running \code{gmsh} would put a combinatorial, non-differentiable step in the middle of
the gradient path, and remeshing between optimizer iterations would make the objective
discontinuous (\at{wheel\_mesh.py}{22--27}). Instead:

\begin{quote}
\textbf{Node positions are a closed-form differentiable function of the 14 genes;
element connectivity is a compile-time constant.}
\end{quote}

\paragraph{Why a clean partition exists at all.}
The spoke is a spiral sweeping about $45^\circ$ of angle inside a $30^\circ$ sector, so
the first guess is that adjacent spokes must overlap. Three measured facts save it
(\at{wheel\_wheel.py}{14--38}):

\begin{enumerate}
\item \textbf{Radius is strictly monotone} along the centerline
  ($12.700 \to 48.900$ in the docstring's own numbers, which predate the rim band being
  thickened inward to \code{RIM\_RADIUS\_MM = 48.5}; the argument is unchanged) while
  $\theta$ rises to $45.16^\circ$ and returns. So each spoke is a single-valued curve
  $\theta(r)$, and twelve rotated copies of a single-valued $\theta(r)$ can never
  intersect. Minimum clearance outside the hub circle: 0.891\,mm.
\item The \textbf{weld footprint} on each ring circle is smaller than the
  sector---$15.85^\circ$ at the hub and $13.18^\circ$ at the rim, of 30
  available---so the twelve junctions tile each ring with a genuine gap.
\item The centerline endpoints are \textbf{locked} on the ring circles, and the end
  cross-section is symmetric about the centerline, so it crosses its ring circle
  exactly at its own midpoint. That corner is available in closed form with \emph{no
  root-find at all}.
\end{enumerate}

Fact 1 makes a structured mesh possible; facts 2 and 3 make the junction blocks
well-shaped instead of slivers.

\paragraph{Seven blocks per sector, twelve sectors}
(\code{BLOCK\_ORDER}, \at{wheel\_wheel.py}{2394}): \code{spoke},
\code{hub\_junction}, \code{rim\_junction}, \code{hub\_collar\_weld},
\code{hub\_collar\_free}, \code{rim\_band\_weld}, \code{rim\_band\_free}. The two ring
blocks per ring are \emph{split at the weld footprint} rather than graded---because
splitting makes the weld arc's two ends into block \emph{corners}, so a contiguous run
of ring nodes coincides with the junction's arc nodes by construction rather than by
arranging a distribution to land on them.

\paragraph{Coons patches.}
The junction blocks are four-sided curvilinear regions, filled by a \emph{bilinearly
blended Coons patch} (\code{coons\_patch}, \at{wheel\_wheel.py}{611}). Given the four
boundary curves,
%
\begin{align}
  X(u,v) \;=\;& \bigl[(1-v)\,\mathrm{bottom}(u) + v\,\mathrm{top}(u)\bigr]
                && \text{ruled in } v \nonumber\\
              &+ \bigl[(1-u)\,\mathrm{left}(v) + u\,\mathrm{right}(v)\bigr]
                && \text{ruled in } u \nonumber\\
              &- \bigl[\text{bilinear interpolation of the four corners}\bigr].
  \label{eq:coons}
\end{align}
%
The subtraction removes the double-counting of the corners; the result interpolates all
four boundaries exactly. This is the classic transfinite interpolation construction.
The code checks corner agreement to $10^{-9}$\,mm on the numpy path, because a swapped
or unreversed edge produces a patch that is \emph{folded} rather than merely ugly, and
the fold surfaces as a negative Jacobian a long way downstream.

\paragraph{Ring stations, and a two-stage root find}
(\code{ring\_station}, \at{wheel\_wheel.py}{296}). Where a flank crosses its ring circle
must be exact, because that point is a corner shared by three blocks. Stage 1 brackets
by finding an actual \emph{sign change} of $r(s) - \mathrm{radius}$ on a dense
grid---deliberately \emph{not} by globally inverting $r(s)$, because an S-shaped
centerline can cross the circle three times and a global \code{interp} lands on the
wrong branch (measured: it left the Coons corner $5\times10^{-3}$ to
$2\times10^{-2}$\,mm out on 2 of 60 feasible genomes). Stage 2 is a \emph{fixed count of
Newton refinements} against the sampler itself, bringing $|x| = \mathrm{radius}$ to
$10^{-15}$. The fixed count is what keeps it traceable \emph{and} differentiable:
unrolled Newton's derivative converges to the implicit-function derivative, so no
custom JVP rule is needed.

\paragraph{Seams and single ownership.}
This is the safety net of the whole module:

\begin{quote}
A seam mismatch produces a mesh that plots correctly, has positive Jacobian everywhere,
and quietly models a wheel with twelve cracks in it. Nothing about the solve complains.
\end{quote}

So every block writes its own coordinates independently, a \emph{union-find} structure
(\code{\_UnionFind}, \code{:2456}, with path compression and lowest-id-wins so ownership
is deterministic) merges the node ids declared equal, and \code{check\_seams} reports
the largest distance between what an owner and a non-owner \emph{independently
computed} for the same node. That number must be at machine precision---measured
$3\times10^{-14}$\,mm for the filleted sector, $7.1\times10^{-15}$ for the tri-block.

\paragraph{Element orientation.}
A polar block indexed $(\theta, r)$ is left-handed, because
$\hat r \times \hat\theta = +\hat z$. Its elements come out with negative Jacobian,
which in an FE assembly is not merely cosmetic---it contributes \emph{negative
stiffness}, which looks to an optimizer like free compliance.
\code{\_orient\_elements} (\code{:2971}) flips whole blocks rather than trusting each
one.

\paragraph{Mesh quality: the scaled Jacobian.}
\code{wheel\_mesh.scaled\_jacobian} (\code{:323}) computes, at each of the four corners
of an element, the cross product of the two edges meeting there normalised by their
lengths---i.e.\ \textbf{the sine of the corner angle}---and takes the worst. A value of
1.0 is a perfect rectangle; near 0 is a collapsed sliver; \textbf{negative is
inverted}. It deliberately measures \emph{shape}, not aspect ratio (a long thin
rectangle scores 1.0), which is the right property for validity. Its
\code{xp=jax.numpy} path is what makes the mesh-validity barrier differentiable.

\subsection{Phase averaging and randomized quasi-Monte Carlo}
\label{sec:phase}

A rolling wheel presents a different spoke pattern to the ground at every angle. The
pattern period is \code{SECTOR\_DEG} $= 30^\circ$. Measured, the axle drop varies
\textbf{7.0\% std/mean and 19.7\% peak-to-peak} over that period---so phase is a design
variable, not a quadrature nuisance, and \code{phase\_ripple} is a real objective term.

\code{phase\_stencil} (\at{wheel\_objective.py}{966}) offers three schemes. The default
is \code{rqmc}: an 8-point uniform stencil shifted by a random offset \emph{drawn from
an 8-point sub-lattice}. This is a quantized randomized-QMC design, and the
quantization is a performance fact rather than statistical taste.

\begin{itemize}
\item \emph{Randomizing at all} matters because the rim's contact faceting lives at
  $0.25^\circ$ while the stencil samples at $3.75^\circ$, so the artefact
  \textbf{aliases}. Under a fixed stencil the aliased artefact is a deterministic
  function of the genes with a nonzero gradient---precisely what an optimizer chases. A
  random shift makes it zero-mean noise instead.
\item \emph{Quantizing} matters because \code{wheel\_wheel.coord\_fn} keys its JIT
  cache on \code{float(phase)}. A \emph{continuously} random offset misses on every
  phase of every step and pays a measured 0.774\,s re-trace eight times per
  step---roughly doubling the actual solving. The union of all reachable phases is a
  fixed $8\times8 = 64$-point lattice, and \code{\_COORD\_FN\_CACHE\_MAX = 128} is sized
  to hold it.
\end{itemize}

\subsection{Projected Adam}
\label{sec:adam}

Stage 3's optimizer is Adam with a cosine learning-rate schedule, global gradient-norm
clipping, and \emph{projection} onto the unit box,
$z \leftarrow \operatorname{clip}(z - \mathrm{step},\,0,\,1)$
(\code{wheel\_stage3.adam\_update}, \code{:248}; \code{project}, \code{:262}).

Projection rather than a sigmoid reparameterisation, for a concrete reason:
\textbf{six of the fourteen genes sit exactly on a bound at the shipped genome}
(\code{cy2} high, \code{t1}/\code{t2}/\code{t3} at \code{MIN\_WALL\_MM}, \code{R\_rim}
high). A squashing reparameterisation has vanishing gradient there, so the run would
begin with six genes frozen. Projection has no such pathology---a gene on a bound whose
gradient points inward moves immediately.

% ===========================================================================
\section{Glossary of repository-specific terminology}
\label{sec:glossary}
% ===========================================================================

\subsection{Structural and process vocabulary}

\begin{termtable}{Structural and process vocabulary.}
\S$N$ / ``section $N$'' &
  A numbered section of \code{PLAN.md}, the project record. Cited constantly in
  comments and commit messages. \S92 is the highest as of 2026-08-29. &
  \code{PLAN.md} \\
\addlinespace
arc &
  An open unit of work with its own plan file (\code{FILLET\_PLAN.md},
  \code{HUBSHARE\_PLAN.md}, \ldots). A ``closed arc'' is finished, its file deleted and
  its record folded into a \code{PLAN.md} section. &
  \at{PLAN.md}{79--115} \\
\addlinespace
M0\ldots M9 &
  Milestones of the original master plan. M2 = mesh validity, M3 = beam agreement,
  M4 = full-wheel linear FEA, M5 = geometric nonlinearity, M6 = real contact,
  M7 = the adjoint, M8a = the objective, M8b = the optimizer, M9 = buckling
  eigenvalue. Sub-milestones such as \code{M8b-i.6} are refinements. &
  \code{PLAN.md}, every study docstring \\
\addlinespace
gate &
  A study driver that \emph{exits nonzero on failure}. Nine of them make up
  \code{make studies}. Distinguished from a \emph{calibration} (measures a constant)
  and a \emph{machine description} (times, not physics). &
  \at{PLAN.md}{1625}\,ff. \\
\addlinespace
G1\ldots G10, S1\ldots S13, A1\ldots A4 &
  Individually numbered checks \emph{inside} a gate. G-numbers are
  \code{study\_gradient}'s, S-numbers \code{study\_stage3}'s, A-numbers
  \code{study\_beam\_agreement}'s. &
  study docstrings \\
\addlinespace
driver &
  A script in \code{studies/} that runs a measurement and writes
  \code{studies/study\_x.json} (and \code{.jpg}). Committed \emph{with} its artifacts,
  by house rule. &
  \code{PLAN.md} header \\
\addlinespace
REDS &
  A closed arc: ``the five inherited reds''---five long-standing failing tests, cleared
  in \S31. &
  \at{PLAN.md}{5095} \\
\addlinespace
red &
  A failing test. \code{make test} reads \code{0 failed} since \S31;
  \code{xfail\_strict = true} means an \code{xfail} that \emph{starts passing} is also a
  failure. &
  \at{pyproject.toml}{31--46} \\
\addlinespace
the shipped genome &
  \code{best\_solution.json}, hash \code{09e8188}, promoted 2026-08-14. Distinct from
  the \emph{GA/beam optimum} \code{36aed36}, which is pinned forever in
  \code{best\_solution\_ga\_beam.json}. &
  \at{PLAN.md}{257--263} \\
\addlinespace
promotion &
  Moving a candidate into \code{best\_solution.json}. Never a one-file change: the
  banner, drivers with measured constants, and \code{make svk} move too;
  \code{tests/test\_promotion.py} carries the checklist. &
  \code{PLAN.md} header \\
\addlinespace
env-opt / env-cad &
  The two virtualenvs. &
  \code{requirements-*.txt} \\
\addlinespace
the hand-off &
  \code{wheel\_fea.py} spawning \code{.venv-cad/bin/python src/wheel\_step\_export.py}
  as a subprocess at the end of a GA run. Deliberately non-fatal. &
  \code{wheel\_fea.py} tail \\
\end{termtable}

\subsection{Geometry vocabulary}

\begin{termtable}{Geometry vocabulary.}
sector &
  One twelfth of the wheel, $30^\circ$. \code{SECTOR\_DEG = 360/12}. The mesh builds
  sector 0 and rotates it eleven times. &
  \at{wheel\_wheel.py}{174} \\
\addlinespace
block &
  One structured quadrilateral grid. Seven per sector unfilleted, eleven filleted,
  twelve for the rim tri-block. &
  \code{BLOCK\_ORDER}, \code{:2394} \\
\addlinespace
seam &
  A shared edge between two blocks, declared as
  \code{(block\_a, side\_a, block\_b, side\_b, dk, reverse)}. \code{dk} is the sector
  offset. &
  \code{\_seam\_table}, \code{:2420} \\
\addlinespace
whole-edge seam &
  A seam where an entire block edge matches an entire other block edge. The
  alternative---a \emph{partial-edge seam}---is avoided by splitting blocks, because
  whole-edge single ownership is ``the whole safety net''. &
  \code{:2404--2419} \\
\addlinespace
span / thick / weld &
  The three element-count directions: \code{n\_span} along the centerline,
  \code{n\_thick} through the thickness, \code{n\_weld} along the weld arc. Three counts
  are \emph{not free} because they sit on shared seams. &
  \code{WheelConfig}, \code{:198} \\
\addlinespace
smoke / coarse / medium / fine &
  Mesh refinement rungs. \textbf{Warning:} \code{wheel\_mesh} and \code{wheel\_wheel}
  both define these names and they are different meshes---a name passed between them
  once resolved silently against the wrong mesh and inflated a convergence order by
  25\%. &
  \code{CONFIGS}, \code{:252}; \at{wheel\_mesh.py}{123} \\
\addlinespace
collar &
  The meshed annulus of the hub, \code{COLLAR\_DEPTH\_MM = 5.0} deep. Inside it the hub
  is treated as a rigid body, because a full disk cannot be one structured quad block
  (a polar grid degenerates at $r = 0$). &
  \code{:195} \\
\addlinespace
rim band &
  The annulus from \code{rim\_inner\_radius()} (48.5) to
  \code{RIM\_OUTER\_RADIUS\_MM} (50.0). 1.5\,mm thick, and it holds about 32\% of the
  wheel's compliance. &
  \code{:177} \\
\addlinespace
junction / weld &
  Where a spoke meets a ring. Two per spoke, 24 in total. &
  \code{sector\_blocks}, \code{:2209} \\
\addlinespace
weld footprint &
  The arc, in degrees, that one junction occupies on its ring circle. &
  \code{weld\_footprints\_deg}, \code{:422} \\
\addlinespace
void / slot &
  \code{SECTOR\_DEG} minus the arc a spoke's material occupies at the hub---the empty
  arc between adjacent spoke roots, into which two fillets must grow from either side.
  Negative means adjacent spokes overlap. &
  \code{hub\_void\_deg}, \code{:467} \\
\addlinespace
arrival angle &
  Angle between centerline and ring \emph{tangent}: $0^\circ$ tangential, $90^\circ$
  radial. See \S\ref{sec:arrival}. &
  \code{arrival\_angles}, \code{:380} \\
\addlinespace
flank &
  One of the two offset curves bounding the spoke. $\eta = +1$ is ``top'', $-1$ is
  ``bottom''. &
  \code{offset\_band} \\
\addlinespace
flank orientation &
  $(\eta_{\text{hub}}, \eta_{\text{rim}})$, each $\pm1$: which flank straddles each
  ring. A \textbf{discrete topological decision}, not a smooth parameter. Only 16 of 60
  feasible genomes share the shipped genome's $(+1, +1)$. Pinned for a whole Stage-3
  run. &
  \code{flank\_orientation}, \code{:559} \\
\addlinespace
straddling flank &
  The one that crosses its ring circle. The other never crosses, which is what makes
  each junction four-sided. &
  \code{junction\_stations}, \code{:594} \\
\addlinespace
ring station &
  The arc-length fraction $s$ at which a flank crosses a ring circle. &
  \code{ring\_station}, \code{:296} \\
\addlinespace
free arc fraction &
  $s_{\text{rim}} - s_{\text{hub}}$: the fraction of centerline that actually flexes
  (67--87\%). \textbf{The wheel's dominant stiffness variable, and the beam model has no
  term for it.} &
  \code{:522} \\
\addlinespace
$P_t$ / $P_c$ &
  The two re-entrant corners of an unfilleted junction. $P_t$ is where the straddling
  flank crosses the ring---\textbf{the corner the real part has}, reproduced to
  0.073\,$\mu$m of arc. $P_c$ is where the mesh's half end cap meets the circle---an
  artefact of the meshing, not a corner of the part. &
  \code{:1013--1030} \\
\addlinespace
uncap &
  Replacing the junction's half end cap with the far flank's own continuation to the
  ring circle---the geometry the exporter's \code{\_embed} actually builds.
  \code{UNCAP\_DEFAULT = (True, 1.0)}; the hub and rim entries do \emph{not} mean the
  same thing. &
  \code{:2115--2206} \\
\addlinespace
embed &
  The exporter's straight tangent extensions burying each spoke end inside its ring, so
  the blunt end caps are absorbed by the boolean union. Its \emph{argmax search} over
  20001 candidates is the piece the mesh deliberately does not reproduce (it would be
  non-differentiable). &
  \code{wheel\_step\_export.\_embed}, \code{:248} \\
\addlinespace
bite &
  Weld penetration in \textbf{root thicknesses}:
  $\mathrm{overlap}/(t^2 \cdot \mathrm{width})$. Derived rather than fitted: a spoke
  burying itself to depth $d$ contributes $t\,d\,W$, so the ratio is exactly $d/t$.
  Replaced a raw mm$^3$ floor that was really a proxy for $t$ and nothing else. Floor
  0.25, honestly labelled a geometric floor, not a calibrated one. &
  \code{wheel\_geometry.junction\_bite}, \code{:395} \\
\addlinespace
fold / fold margin &
  See \S\ref{sec:fold}. &
  \code{self\_intersection\_margin}, \code{:421} \\
\addlinespace
fillet\_a / fillet\_b &
  The two blocks per junction of the eleven-block filleted sector. \code{a} is the
  boundary layer along the fillet arc; \code{b} is the wedge crossing the ring circle.
  They are split at $N$ because one edge with two partners would be a partial-edge
  seam. &
  \code{:986--1010} \\
\addlinespace
layer profile &
  The pair \code{(entry, end)} controlling the filleted boundary layer's width
  profile---a Hermite cubic. Since \S85 the default is a \emph{per-genome rule}, not a
  global pair. &
  \code{:1071--1140} \\
\addlinespace
cliff &
  The value of \code{entry} at which a junction's layer reaches zero width. \S88 found
  it has a \textbf{closed form}, $\max_u (Z - a(u))/b(u)$, agreeing with the 90-halving
  bisection it replaced to 2\,ulp. &
  \code{\_layer\_cliff\_from\_scalars}, \code{:1537} \\
\addlinespace
sector-fit clamp &
  Pulls a gene-derived fillet radius back inside the room its own sector has, at 95\% of
  the limit. Took a drawn genome box from 10/16 building to 16/16. A clamp is \emph{not}
  a gate: it keeps the genome and models a \emph{smaller} fillet, which is honest only
  because \code{mesh.fillet\_radii\_mm} and \code{mesh.fillet\_clamped} report what was
  actually built. &
  \code{SECTOR\_FIT\_CLAMP}, \code{:1171} \\
\addlinespace
tri-block &
  The three-quad Y-partition of the rim junction at the faithful (\code{uncap} blend
  0.0) rim, where the region becomes a curvilinear \emph{triangle}. Built and measured
  (\S53), not adopted. &
  \code{studies/study\_tri\_block.py} \\
\end{termtable}

\subsection{Optimization and physics vocabulary}

\begin{termtable}{Optimization and physics vocabulary.}
T1 / T2 / T3 &
  The three cost tiers of the objective. \textbf{T1} is gene-space
  (microseconds, no mesh), \textbf{T2} mesh-space (milliseconds,
  \code{mesh\_coords} only), \textbf{T3} solve-space ($\sim0.7$\,s per phase). The
  tiering is economics, not organisation: a line search can re-evaluate T1 alone. &
  \at{wheel\_objective.py}{30--42} \\
\addlinespace
barrier vs.\ objective term &
  See \S\ref{sec:barriers}. \code{BARRIER\_TERMS} / \code{OBJECTIVE\_TERMS}, with an
  assert that every term is classified. &
  \code{:394--407} \\
\addlinespace
QoI &
  Quantity of Interest---a scalar
  $Q(\code{coords}, \code{u\_full}, \code{y\_ground})$ the adjoint differentiates. The
  \code{QOI} registry holds \code{contact\_force}, \code{strain\_energy}, \code{mass},
  \code{pnorm\_stress}. &
  \at{wheel\_adjoint.py}{450} \\
\addlinespace
axle drop &
  The wheel's deflection under load, in mm. Target 2.0. Measured as the prescribed
  ground indentation at which the contact force equals the service load. &
  \code{solve\_wheel\_contact}, \at{wheel\_fem.py}{1867} \\
\addlinespace
phase / \code{phase\_deg} &
  The rolling angle of the wheel under a stationary ground. Lives on
  \code{build\_wheel(phase\_deg=...)}, \emph{not} on the load---the ground does not
  move. Putting phase in the load pushes the wheel sideways, a different load case, and
  silently breaks the 12-fold periodicity check. &
  \code{\_sector\_coords}, \code{:2583} \\
\addlinespace
stencil &
  The set of phases one objective evaluation averages over. See \S\ref{sec:phase}. &
  \code{phase\_stencil}, \code{:966} \\
\addlinespace
utilisation (\code{util}) &
  $K_t\,\sigma_{\text{nominal}} / \mathrm{allowable}$. 1.0 is the wall. Shipped genome:
  0.8201 (hub), 0.5454 (rim). &
  \code{t3\_terms}, \code{:1045} \\
\addlinespace
the knee &
  \code{MARGIN\_KNEE\_UTIL = 0.80}. \code{stress\_margin} is
  $\mathrm{soft\_barrier}(\code{util} - 0.80)$---the same function as the \code{stress}
  wall with the knee moved in. Encodes the policy ``margin below 0.8 is close to
  worthless, above 0.95 close to priceless''. Explicitly labelled a judgement, not a
  measurement. &
  \code{:343} \\
\addlinespace
\code{min\_sj} &
  The mesh-validity barrier, summed over \emph{every} element rather than applied to
  the worst---because \code{min} is a max-like kink whose gradient is a delta on
  whichever element is currently worst, and because summing also tells the optimizer
  \emph{how many} elements are marginal. Target 0.2. &
  \code{t2\_vector}, \code{:931} \\
\addlinespace
coupling term &
  The $(\partial Q/\partial\delta)(d\delta^*/dp)$ half of the total gradient at constant
  force. See \S\ref{sec:adjoint}. &
  \code{service\_qoi\_value\_and\_grad} \\
\addlinespace
frozen / freezing &
  Holding a discrete decision fixed so a derivative exists: flank orientation, seam
  ownership, fillet roots, the clamp's decision. Always \emph{named}, never hidden. &
  \code{mesh\_coords} docstring \\
\addlinespace
inert term &
  A term with a large loss value and a tiny gradient---dominates the table, contributes
  nothing to descent. \code{400*n\_inflections} was this in pure form. Gate 8 exists to
  detect it. &
  \code{studies/study\_objective.py} \\
\addlinespace
event &
  A structured record in a run log of something the optimizer decided:
  \code{t1\_reject}, \code{orientation\_flip}, \code{fidelity\_check\_failed},
  abandonment. &
  \code{descend}, \code{:410} \\
\addlinespace
selection key &
  The three-tier promotion sort. See \S\ref{sec:selectionkey}. &
  \code{:195} \\
\addlinespace
\code{make <target>} &
  Every measurement in this tree has a Makefile target: \code{make fillet},
  \code{filletblock}, \code{corner}, \code{hubcap}, \code{gci}, \code{triblock},
  \code{svk}, \code{studies}, \code{test}. &
  \code{Makefile} \\
\end{termtable}

% ===========================================================================
\section{Architecture and module breakdown}
\label{sec:arch}
% ===========================================================================

\begin{lstlisting}
src/
  project_paths.py     35 lines   ROOT/SRC/STUDIES/EXPORT.  STDLIB ONLY -- imported by
                                  wheel_fea, which the CAD env imports.
  jax_config.py        36        `jax.config.update("jax_enable_x64", True)`, and raises
                                  loudly if something traced first.  float64 is mandatory:
                                  at float32 a Newton solve cannot reach 1e-10, an FD check
                                  has no plateau, and the patch test's 1e-12 is below eps.
  wheel_genome.py     183        names, bounds, normalisation, genome_hash, record I/O.
                                  numpy+stdlib only.
  wheel_geometry.py   436        THE GEOMETRY KERNEL.  Bezier, taper, offset band,
                                  curvature, fold margin, junction bite.  Dual-backend via
                                  `xp=`; never imports jax.
  wheel_mesh.py       391        One spoke block: sampler, grid connectivity (Q4/Q9),
                                  boundary node sets, quality metrics (scaled Jacobian,
                                  aspect ratio, areas).
  wheel_fea.py       1713        STAGE 1/2.  Constants, GENE_SPACE, Castigliano beam model,
                                  the nine loss terms, the pygad driver, the fillet-arc
                                  plotting helpers, and the CAD hand-off.  numpy only.
  wheel_wheel.py     3709        THE FULL 360 DEGREE MESH.  Ring geometry, block builders,
                                  Coons patches, the filleted eleven-block sector, seam
                                  tables, union-find assembly, WheelMesh, the differentiable
                                  `mesh_coords` + jitted `coord_fn`, area and quality
                                  reports.  Largest and most intricate file in the tree.
  wheel_fem.py       1914        THE FE KERNEL.  Shape functions, energy, autodiff'd force
                                  and tangent, DofMap, penalty contact, linear and Newton
                                  solvers, stress recovery, spoke/wheel/contact problems.
  wheel_adjoint.py    759        THE GRADIENT.  QoIs, the adjoint solve, load-controlled
                                  total derivatives, LOBPCG buckling eigenpair.
  wheel_objective.py 1461        THE STAGE-3 SCALAR.  Fourteen terms in three tiers, the Kt
                                  stress constraint, the hub fillet cap model, phase
                                  stencils, jit caches, loss breakdown reporting.
  wheel_stage3.py    1093        THE DRIVER.  Projected Adam (and an L-BFGS-B alternative),
                                  step-reject policy, selection key, run records, CLI.
  wheel_pool.py       408        Parent half of the phase worker pool.  stdlib+numpy, NO jax.
  wheel_pool_worker.py 103       Child.  Pins threads BEFORE its first import.
  wheel_step_export.py 1248      THE EXPORTER (CadQuery env).  Profile construction,
                                  embedding, boolean union, fillet ladder, health checks,
                                  STEP write, manifest.
studies/              ~19 drivers + their committed .json/.jpg artifacts
tests/                25 files, 503 test functions
\end{lstlisting}

\subsection{The import graph, and the constraint that shapes it}

\begin{lstlisting}
                 project_paths  (stdlib)
                       |
      wheel_geometry --+-- wheel_fea ------> wheel_step_export   [CAD ENV]
       (xp-generic)         (numpy)              (cadquery/OCP)
            |                 |
        wheel_mesh            |
            |                 |
            +---> wheel_wheel + jax_config
                       |
                  wheel_fem
                       |
                 wheel_adjoint
                       |
                wheel_objective
                       |
                 wheel_stage3 ---- wheel_pool ---- wheel_pool_worker
\end{lstlisting}

The horizontal line matters more than the vertical one. Everything to the left of and
including \code{wheel\_fea} must import \textbf{with numpy alone}, because
\code{wheel\_step\_export.py} runs in the CadQuery interpreter and imports
\code{wheel\_fea} and \code{wheel\_geometry} from it. A single top-level
\code{import jax} anywhere in that half breaks the exporter, and
\code{tests/test\_import\_hygiene.py} runs \code{import wheel\_fea} in a jax-free
subprocess to prove it has not happened.

This is also why \code{wheel\_wheel.mesh\_coords} imports \code{jax.numpy}
\emph{lazily}, inside the function, and why \code{wheel\_geometry} takes \code{xp} as a
parameter instead of importing an array module. Note also that
\code{MAX\_ARRIVAL\_DEG}, \code{MIN\_FOLD\_MARGIN\_MM} and \code{MIN\_JUNCTION\_BITE}
live in \code{wheel\_geometry} rather than where they are used, because their
\emph{producer} and \emph{consumer} are on opposite sides of the interpreter boundary
and this is the only module both can import.

\code{src/} reaches the interpreter three ways---\code{pyproject.toml}'s
\code{pythonpath} for pytest, \code{export PYTHONPATH} in the Makefile, and the root
\code{conftest.py} (which also seeds \code{os.environ} so the three subprocess-spawning
tests behave identically under bare \code{pytest} and under \code{make test}). A package
with \code{\_\_init\_\_.py} was rejected precisely because an \code{\_\_init\_\_}
importing jax-dependent siblings would break the CAD environment.

\subsection{Caches, and why they are keyed the way they are}
\label{sec:caches}

There are three JIT caches in this tree, all following the same idiom and all with the
same failure mode documented.

\paragraph{\code{wheel\_wheel.\_COORD\_FN\_CACHE}} (\code{:2760}), the jitted
$\mathrm{genes} \to \mathrm{coords}$ map. The JIT is load-bearing, not tuning:
\code{sector\_blocks} unrolls a fixed-count Newton over \code{n\_curve} stations for
every ring crossing, so tracing is expensive, and \code{jax.vjp} on an \emph{untraced}
closure re-traces on \textbf{every call}---measured at 0.774\,s against 0.05\,s for the
entire rest of the adjoint. Without the cache, 97\% of a Stage-3 gradient is re-deriving
a jaxpr that never changes.

The cache \textbf{cannot be keyed on the mesh object}---the obvious thing, and wrong.
Every finite difference, every sweep point and every optimizer step builds a \emph{new}
mesh at new genes, so a per-object cache misses on literally every call it exists to
serve while looking like it works. The key is the \emph{static recipe} instead: element
counts, orientation, ownership bytes, phase, \code{repr(uncap)}, and the fillet recipe.
Genes are the traced \emph{argument} and must never be in the key. Correspondingly, the
filleted mesh's frozen roots are passed as a \emph{traced array argument} rather than
closed over---because they depend on the genome, and closing over them would put a
design into the jaxpr and re-trace on every step (2.78\,s to trace once, 0.0095\,s per
call after; 2.78\,s \emph{every} call if closed over).

\code{repr(uncap)} in the key is a correctness entry, not a performance one. Its absence
was invisible while \code{UNCAP\_DEFAULT} was \code{False}; when \S38 flipped the
default, a mesh built with an explicit \code{uncap=False} got traced coordinates from
the \emph{other} geometry---\textbf{0.448\,mm of error} against the $10^{-9}$\,mm the
path is documented to hold to.

\paragraph{\code{wheel\_objective.\_T1\_CACHE}} (\code{:872}) and
\textbf{\code{\_KT\_CACHE}} (\code{:521}) follow the same pattern for the T1 term vector
and the $K_t$ value-and-gradient. T1's jit took a full evaluation's \code{t1\_vector}
from about 1.06\,s to 0.0054\,s---a roughly $195\times$ improvement that mattered
disproportionately, because once the phase loop was parallel, T1 and T2 (which run in
the parent on every path) became the whole of Amdahl's serial fraction.

\subsection{The phase pool}
\label{sec:pool}

\code{wheel\_pool.py} runs the eight per-phase solves of one evaluation in eight
processes. It is deliberately \emph{not} \code{multiprocessing.Pool}, for three
load-bearing reasons.

\begin{enumerate}
\item \textbf{Thread settings must be in place before numpy and JAX are imported.}
  \code{OMP\_NUM\_THREADS} and its siblings are read by OpenBLAS/MKL at import and
  \code{XLA\_FLAGS} by XLA at its first trace; \code{spawn}'s \code{initializer=} fires
  \emph{after} the unpickling that already imported everything. \code{fork} would be
  early enough and is unusable on two counts (forking an initialised JAX process is a
  documented hazard; \code{fork} does not exist on Windows). So the pool uses an
  explicit \code{Popen(env=...)}.
\item \textbf{Phase slots must pin to workers}, so slot $i$ always goes to worker
  $i \bmod n_{\text{workers}}$ and each worker only ever traces its own phases' jaxprs.
\item \textbf{Results must combine in slot order.} Floating-point addition is not
  associative, so an as-completed reduction would make the aggregate depend on the
  scheduler. In slot order the arithmetic is identical, which is what lets S13 gate
  \emph{pooled $=$ serial exactly} rather than to a tolerance.
\end{enumerate}

\code{XLA\_FLAGS} deserves its own note, because it is a \textbf{correctness} setting
here. Measured: two plain serial runs of one \code{coarse} adjoint, in two separate
interpreters with no pool anywhere, agree on every forward value to the bit and
disagree on the \emph{gradient} by $3.33\times10^{-16}$---XLA's CPU intra-op thread pool
is sized from the machine and its parallel reductions do not associate the same way
twice. Pinned to one thread, the same comparison is exactly zero, and it costs 19.84\,s
against 20.43\,s. Nothing in OMP/MKL/OPENBLAS/NUMEXPR touches that pool.

Nothing unpicklable crosses the boundary: \code{WheelMesh.\_\_slots\_\_} holds
\code{\_coord\_fn}, a jitted \code{PjitFunction} after the first \code{mesh\_coords}
call. So the worker builds its own mesh from the genes and returns only the leaves
\code{t3\_terms} reads---three floats and three 14-vectors. The reply is shaped
\emph{exactly like} the serial adjoint dictionary, so the loop body in
\code{t3\_terms} is literally the same lines on both paths.

Measured: $3.95\times$ at 8 workers on 16 cores; a production run projects 46.46\,h down
to 11.77\,h.

% ===========================================================================
\section{Full workflow walkthrough}
\label{sec:workflow}
% ===========================================================================

Let me trace one concrete path end to end: \textbf{Stage 3 takes one Adam step} on the
shipped genome. This is the deepest path in the system and touches almost everything.

\subsection*{Step 0 --- the CLI}

\begin{lstlisting}[language=bash]
.venv-opt/bin/python src/wheel_stage3.py --start best --steps 100 \
    --config medium --kinematics svk --min-wall 1.2 --workers -1
\end{lstlisting}

\code{main()} (\at{wheel\_stage3.py}{932}) loads \code{best\_solution.json} via
\code{load\_genes} (\code{:844}), converts the 14-key dictionary to an ordered vector
with \code{wheel\_genome.genes\_to\_vector}, and normalises it into the unit box.
\code{--min-wall 1.2} calls \code{wheel\_fea.set\_min\_wall}, which rebuilds
\code{GENE\_SPACE}'s thickness bounds---this is \emph{not} cosmetic: \code{descend}
projects its start into the box before it steps, so loading the shipped 1.2\,mm genome
without this flag would silently lift all four thicknesses to 2.0 and optimise a
different, heavier wheel without a word of complaint
(\at{wheel\_fea.py}{216--236}).

\subsection*{Step 1 --- pinning the topology}

\code{descend} (\code{:410}) computes the \textbf{flank orientation} once, from the
starting genes:

\begin{lstlisting}[style=py]
orientation = WW.flank_orientation(wg.denormalize(z, low, high), wcfg, span_mm=span_mm)
\end{lstlisting}

\code{flank\_orientation} (\at{wheel\_wheel.py}{559}) samples the two flanks at the
endpoints and asks which one is inside its ring circle. The answer is a pair of $\pm1$.
It is \textbf{pinned for the whole run} and threaded through every \code{build\_wheel}
call thereafter, because it is a discrete topological decision: a step that flips it
changes which block owns which node, and the gradient of that step does not exist. It is
\emph{also} re-derived every step for reporting, and a disagreement is logged as an
\code{orientation\_flip} event---the run keeps going on the pinned topology, and the
event says the trajectory has walked somewhere the frozen answer is stale.

Then the pool is created (\code{WP.PhasePool(n\_workers)}), owned by \code{descend}
rather than by the \code{Evaluator}, so the \code{finally} that closes it sits next to
the loop a Ctrl-C can interrupt.

\subsection*{Step 2 --- drawing the phase stencil}

\begin{lstlisting}[style=py]
phases = WO.phase_stencil(n_phase=8, n_sub=8, scheme="rqmc", rng=rng)
\end{lstlisting}

Eight angles in $[0^\circ, 30^\circ)$, on the fixed 64-point lattice
(\S\ref{sec:phase}).

\subsection*{Step 3 --- building eight meshes}

\code{Evaluator.\_\_call\_\_} (\code{:377}) calls
\code{WO.phase\_meshes(genes, cfg, phases, orientation=orientation)}, which is
twelve-fold \code{build\_wheel(genes, cfg, phase\_deg=p, orientation=orientation)}.
(Pooled, the parent builds only mesh 0---T2 reads \code{meshes[0]}---and each worker
builds its own.)

Inside \code{build\_wheel} (\at{wheel\_wheel.py}{2869}):

\begin{enumerate}
\item \code{\_sector\_coords} (\code{:2583}) calls \code{sector\_blocks} (\code{:2209}),
  which builds sector 0's seven node grids.
  \begin{itemize}
  \item \code{global\_sampler} (\code{:275}) wraps \code{wheel\_mesh.band\_sampler} and
    shifts it into the global frame, so \emph{every} block that touches the band goes
    through the same closure and shared nodes agree bitwise rather than to $O(h^2)$.
  \item \code{junction\_stations} then \code{ring\_station} (\S\ref{sec:meshing}) finds
    $s_{\text{hub}}$ and $s_{\text{rim}}$ to $10^{-15}$.
  \item the \textbf{spoke} block is \code{sample(s\_grid[:,None], eta\_grid[None,:])}
    --- the offset band of Equation~\eqref{eq:offsetband}, sampled uniformly in arc
    length $\times$ $\eta$.
  \item each \textbf{junction} is a Coons patch whose four sides are: the spoke's own
    end cross-section (the same array, so the seam is exact by construction), an arc
    along the ring circle, the far flank sampled at $-\eta$, and a straight segment to
    the far end. Corner angles come out $79.5 / 90 / 90 / 79.5$ degrees rather than the
    10.5 the near-tangent arrival suggests---that is the \S\ref{sec:arrival} inversion
    paying off.
  \item each \textbf{ring} is split into a \code{\_weld} polar block (spanning exactly
    the weld arc) and a \code{\_free} polar block (the rest of the sector), both laid
    out in increasing $\theta$ so that
    \code{weld.i1} $\to$ \code{free.i0} $\to$ next \code{weld.i0} holds for every genome
    and the tiling closes.
  \end{itemize}
\item The seven grids are rotated into twelve sectors by \code{\_rotate} and
  concatenated. \textbf{Phase is applied here}, not to the load.
\item The \textbf{seam table} is walked; a \code{\_UnionFind} merges declared-equal
  nodes; the seam error is measured \emph{before} the non-owners' coordinates are
  discarded.
\item Nodes are relabelled to a compact global numbering; \code{owners} (the raw index
  of each merged node) is saved on the mesh---this is what \code{mesh\_coords} will
  index with.
\item Connectivity is emitted per block by \code{wheel\_mesh.grid\_connectivity} and
  passed through \code{\_orient\_elements} to flip left-handed polar blocks.
\item \code{\_node\_sets} / \code{\_edge\_sets} build \code{hub\_tie},
  \code{rim\_outer} and friends. Edge \emph{segments}, not just node sets, because a
  distributed load needs consistent nodal forces and the correct weights on a quadratic
  edge are $1/6, 4/6, 1/6$---not equal thirds.
\end{enumerate}

Out comes a \code{WheelMesh} carrying coordinates, connectivity, config, per-element
block and region tags, node and edge sets, \code{seam\_error\_mm}, \code{owners},
\code{orientation}, \code{phase\_deg}, \code{uncap}, and (if filleted) the applied radii
and the frozen root record. At \code{medium} this is about 104k degrees of freedom.

\subsection*{Step 4 --- the objective}

\code{WO.objective(z, cfg, normalized=True, ...)} (\at{wheel\_objective.py}{1308})
denormalises, then runs three tiers.

\paragraph{T1} (\code{t1\_vector}, \code{:774})---genes only, no mesh, microseconds via
the jitted cache. Seven terms, returned as a \emph{vector} so \code{jax.jacrev} gives
every term's gradient in one backward pass (which is what the per-term inertness census
needs):

\begin{itemize}
\item \code{x\_order}---control points must stay ordered with a 2\,mm gap, plus a
  fold-back barrier on $-\min(\mathrm{diff}(x))$;
\item \code{hub\_overlap}---root thickness plus two fillet radii plus clearances must
  fit the chord $2 R_{\text{hub}}\sin(\pi/12) \approx 6.57$\,mm available to one sector;
\item \code{smoothness}---\textbf{rewritten} for differentiability. The old
  \code{400*n\_inflections} was an integer from \code{count\_nonzero} behind a
  data-dependent mask with \emph{exactly zero gradient} while being 20.6\% of the GA's
  loss. It is replaced by a curvature-rate integral $\int (d\kappa/ds)^2\,ds$ plus a
  smooth sign-reversal product
  $\sum \max\bigl(0, -\kappa_i \kappa_{i+1}/\kappa_{\text{ref}}^2\bigr)\,\Delta s$;
\item \code{fold}---\S\ref{sec:fold};
\item \code{arrival}---\S\ref{sec:arrival};
\item \code{fillet}---a closed-form geometric feasibility margin for each of four fillet
  corners, barriered individually so one infeasible fillet cannot be averaged away by
  three comfortable ones. \textbf{This is what gives \code{R\_hub}/\code{R\_rim} a live
  gradient} even though the default mesh models no fillets;
\item \code{fillet\_cap}---\code{R\_hub} may not exceed \code{hub\_fillet\_cap\_mm}, the
  largest hub fillet the part can actually build (\S\ref{sec:hubcap}). A separate term
  from \code{fillet}, deliberately: \code{fillet} asks whether a circle fits in one
  re-entrant \emph{corner}, this asks whether it fits in the \emph{slot} two spokes
  leave between them. Different mechanism, different gradient, and every census in the
  tree reads at term granularity---folding them together would make ``is the cap
  binding?'' unanswerable from the loss table.
\end{itemize}

Plus \code{buckling}, computed from the numpy beam surrogate with an explicitly
asserted \emph{zero gradient}---currently exactly 0.0 (ratio 0.139 against a threshold
of 1.0), and the gate asserts it \emph{stays} zero, so the day it starts to bite the
census fails and says so.

\paragraph{T2} (\code{t2\_vector}, \code{:931})---mesh-space, via \code{mesh\_coords}
alone:

\begin{lstlisting}[style=py]
coords = WW.mesh_coords(genes, mesh)          # differentiable, jitted
mass_g = sum(element areas) * width * density
min_sj = sum_e soft_barrier(0.2 - sj_e)
\end{lstlisting}

\paragraph{T3} (\code{t3\_terms}, \code{:1045})---one
\code{service\_qoi\_value\_and\_grad} per phase (\S\ref{sec:adjoint}), each of which is
a secant to service force, a re-solve at the found indentation, one stiffness assembly,
one factorisation, and back-solves for \code{contact\_force}, \code{pnorm\_stress} and
whatever else was requested, all sharing one mesh vjp. Then:

\begin{itemize}
\item $\code{deflection} = 2500 \left(\frac{\overline{\mathrm{drop}} - 2.0}{2.0}\right)^2$,
  the mean over the stencil;
\item \code{stress}---a smoothed maximum over the \emph{phase} samples ($p$-norm at
  $p = 8$, replacing CVaR, which at 8 samples is ``the worst one or two'' and whose
  gradient jumps as the set switches) fed into the
  $K_t\,\sigma_{\text{nom}} \le \mathrm{allowable}$ barriers, one per junction
  (\S\ref{sec:stress});
\item \code{stress\_margin}---the same two utilisations, priced instead of walled, with
  the knee at 0.80;
\item \code{phase\_ripple}---std/mean of the drops over the stencil.
\end{itemize}

Note that the two exponents in T3 are \textbf{not the same quantity}:
\code{stress\_phase\_p = 8.0} aggregates the eight phase samples;
\code{stress\_gauss\_p = 4.0} is the volume-weighted $p$-norm exponent over the Gauss
points inside one solve.

\code{objective} sums the fourteen values, sums the fourteen gradients, applies the
unit-box chain rule if \code{normalized}, and returns \code{(value, grad, breakdown)}
where the breakdown carries per-term value, share, gradient norm and gradient share.

\subsection*{Step 5 --- the step}

Back in \code{descend}:

\begin{lstlisting}[style=py]
lr_t             = cosine_lr(lr, i-1, steps)
g_clipped, g_norm = clip_global_norm(grad, 1.0)
delta, m, v      = adam_update(g_clipped, m, v, i, lr_t)
z_trial          = project(z - scale*delta)
\end{lstlisting}

Then the \textbf{T1 pre-check}: \code{t1\_barrier\_sum(z\_trial, ...)}---microseconds,
no mesh, no solve---and if the trial's barrier sum exceeds
$\max(\code{T1\_REJECT}, b_{\text{here}})$ it is refused \emph{before any mesh is
built}.

\code{T1\_REJECT = 1.0e4}, and the number is the whole lesson
(\at{wheel\_stage3.py}{143--167}). It was 1.0, with the rule ``never step further into
the wall''. That \textbf{deadlocked} a probe: the shipped design carries
$\code{hub\_overlap} = 52.13$, the stress-reducing direction raises it to 55.41, and
halving only walks it back to 52.54, so every trial at every scale was refused---194
events, zero accepted steps, thirty-five minutes standing still. The error was
conceptual: every barrier is \emph{already} a term in the objective with a weight
chosen to price it, so a pre-check that also forbids increasing it duplicates and
overrides that job---and what it overrides is exactly the constrained trade the run
exists to make. The screen now keeps only the job the objective cannot do: refusing to
spend a solve on geometry that will not mesh. Healthy barriers are tens; an analytically
infeasible fillet scored 112{,}000.

If the objective \emph{raises}---\code{NewtonDivergedError}, or the
$dF/d\delta \le 0$ guard---the step is halved and retried up to \code{--max-rejects}. On
exhaustion the iterate is restored and the learning rate decays.
\code{wheel\_objective.py} contains \textbf{no \code{try}/\code{except} anywhere},
deliberately: the caller owns the policy, and the policy is in \code{wheel\_stage3}.
This matters beyond robustness, because $\mathrm{slope} \ge 0$ in the Newton loop is the
tangent losing positive definiteness, and swallowing it as a large loss would let the
optimizer walk into a buckled design and score it.

If accepted, the \textbf{warm vector} for the next step is the per-phase indentations,
read straight off \code{breakdown["report"]["rows"]}---because
\code{service\_qoi\_value\_and\_grad}'s \code{delta0} \emph{is} the number reported as
\code{axle\_drop\_mm}. Between two Adam steps the design barely moves, so the previous
answer is a far better guess than a cold linear solve.

Every $N$ steps, an optional \textbf{fidelity check} forward-evaluates the just-accepted
iterate at a second config (\code{medium}), T3 tier only, and \emph{discards the
gradient}---proven by a test that runs the same seed with the feature on and off and
asserts \code{z}, \code{grad} and \code{loss} are bit-identical at every step. Pure
observation: it can report on the trajectory, never redirect it.

\subsection*{Step 6 --- selection and persistence}

Every step is scored by \code{selection\_key} (\S\ref{sec:selectionkey}) and the best is
tracked. The run record (\code{\_record}, \code{:665}) is written incrementally, so a
killed run leaves a readable partial. At the end, \code{--best-out} writes a genome
record through \code{wheel\_genome.save\_record}---genes plus top-level \code{search},
\code{loss\_terms} and \code{metrics} blocks.

\subsection*{Step 7 --- promotion and export}

A promotion moves the candidate to \code{best\_solution.json}, and by house rule that is
\textbf{one atomic commit that is never one file}: the banner, the drivers carrying
measured constants, and \code{make svk} move with it, and
\code{tests/test\_promotion.py} carries the checklist.

Then \code{make export} runs \code{wheel\_step\_export.py} in the CAD environment:

\begin{enumerate}
\item \code{load\_genome} reads \code{best\_solution.json}; \code{warn\_if\_stale}
  compares mtimes.
\item \code{spoke\_edges\_global} re-derives the two flanks from
  \code{wheel\_fea.thicken\_3taper\_curve}---the \emph{same} geometry kernel the mesh
  uses.
\item \code{\_embed} (\code{:248}) appends \textbf{straight tangent segments} to each
  flank end, burying the spoke root 0.5\,mm inside the hub disk and running the tip
  0.25\,mm past $\varnothing$100. Straight, because the previous approach---moving the
  last point and interpolating a spline through it---put a \emph{self-intersecting loop}
  in every spoke: the first span jumped 3\,mm while its clamped tangent pointed
  $40.7^\circ$ against a local curve direction of $76.1^\circ$. The union trimmed the
  loops away but left a 0.041--0.067\,mm near-cusp running the full 22.4\,mm as a knife
  edge. OCC calls that valid (\code{BRepCheck\_Analyzer} does not test for cusps) and
  Parasolid refuses to build a face across it, drops the face, and \textbf{the spoke
  imports open}. A straight line adds no curvature, cannot cusp, and continues the
  flank's own end tangent so the join is $G^1$.
\item \code{build\_profile} unions hub disk $+$ 12 spokes $+$ rim band into one planar
  face, extrudes it, and clips to $\varnothing$100.
\item \code{\_junction\_edges} finds the full-height vertical edges and measures each
  one's \emph{material wedge angle} by classifying sample points against the solid. A
  corner is worth filleting when its wedge exceeds $180^\circ$ (re-entrant). The five
  families separate with a wide margin: $354^\circ$ (the spoke-to-spoke notch at the
  hub, fillet), $152^\circ$ (convex, skip), $180^\circ$ (straight through, no corner),
  $194$--$356^\circ$ (the rim junctions, fillet), $178^\circ$ (the OD trim seam, skip).
\item \code{fillet\_junctions} walks a \emph{ladder}: request the genome's radius; if
  OCC refuses, retry at $0.85\times$ (\code{FILLET\_LADDER\_DECAY}), down to
  \code{MIN\_CURVATURE\_RADIUS\_MM = 0.25}. All 48 corners (24 hub $+$ 24 rim) are
  filleted on the shipped genome, at the requested radii, with \code{kt\_error} 0.0\% at
  both junctions.
\item \code{step\_health} runs a surface census, a minimum-curvature scan (a fault
  detector: 0.25\,mm is well under any fillet requested, so a violation always means a
  construction fault), and \code{despecialize} (converting swept and offset surfaces to
  canonical ones, which is what some importers need).
\item \code{\_export\_with\_settings} writes AP214 STEP; \code{report} writes
  \code{export/wheel\_step\_manifest.json} with the genome hash, junction overlaps
  \emph{and bites}, the fillet radii OCC actually managed versus what the optimizer
  priced, and the solid volume with and without fillets.
\end{enumerate}

For the shipped genome: 39224.5\,mm$^3$ solid, 36145.8\,mm$^3$ unfilleted,
\textbf{3078.77\,mm$^3$ of fillet $=$ 7.85\% of the solid}, 48.64\,g of PLA.

% ===========================================================================
\section{Key algorithms, in detail}
\label{sec:algorithms}
% ===========================================================================

\subsection{\texttt{mesh\_coords} --- the differentiable mesh}
\label{sec:meshcoords}

\code{wheel\_wheel.mesh\_coords(genes, mesh)} (\code{:2633}) returns the mesh's node
coordinates as a differentiable function of the genes. The trick is what it
\emph{freezes}:

\begin{itemize}
\item the \textbf{flank orientation}---taken from \code{mesh}, not recomputed;
\item the \textbf{seam ownership}---\code{coords\_all[mesh.owners]};
\item on a filleted mesh, the tangency and ring-crossing \textbf{roots}, seeded from
  the eager build and refined by one Newton step;
\item the four geometric refusals \code{\_fillet\_curves} can make.
\end{itemize}

So this is \textbf{the derivative of a fixed mesh whose nodes move}, which the docstring
argues is the only derivative that means anything: a step that flips the orientation has
no finite-difference plateau, and \emph{should not}---it is a real discontinuity of the
design space, not a defect in the gradient.

Crucially, \code{build\_wheel} itself is deliberately \emph{not} made traceable. Its
seam check, element orientation pass and validity guards all need concrete coordinates,
and they are the reason the mesh can be trusted; wrapping them in a ``skip when
tracing'' branch is exactly how guards stop running. Instead the traced half is factored
out (\code{\_sector\_coords}) and shared, and \code{mesh\_coords} \emph{reproduces} the
eager result---checked to $10^{-9}$\,mm by \code{study\_gradient.py} gate 3 and
\code{tests/test\_gradient.py}.

\subsection{The \texorpdfstring{$p$}{p}-norm family, and why \texttt{max} never appears in a gradient}
\label{sec:pnormfamily}

Three separate places replace a \code{max} or \code{min} with a smooth aggregate, and
each records the same argument in its own terms.

\begin{table}[h]
\centering
\small
\begin{tabularx}{\linewidth}{@{}>{\raggedright\arraybackslash}p{0.22\linewidth}
                               >{\raggedright\arraybackslash}p{0.10\linewidth}
                               >{\raggedright\arraybackslash}p{0.24\linewidth}
                               X@{}}
\toprule
\textbf{Quantity} & \textbf{Naive} & \textbf{Used instead} & \textbf{Why} \\
\midrule
peak stress over Gauss points & \code{max} & volume-weighted $p$-norm, $p = 4$ &
  argmax jumps as the contact patch sweeps; also \code{max} diverges under refinement \\
\addlinespace
peak over phase samples & CVaR & $p$-norm, $p = 8$ &
  at 8 samples CVaR is ``the worst one or two'' and its gradient jumps as the set
  switches \\
\addlinespace
worst element quality & \code{min sj} &
  $\sum_e \mathrm{soft\_barrier}(0.2 - sj_e)$ &
  the gradient of \code{min} is a delta on whichever element is worst; and the sum tells
  you \emph{how many} are marginal \\
\addlinespace
effective fillet radius & $\min(R, \mathrm{cap})$ &
  \code{smooth\_min(R, cap, k)} &
  needs $C^1$ at the tie, but \textbf{exact} outside a blend of width $k$ \\
\bottomrule
\end{tabularx}
\caption{Four places a hard extremum was replaced, and the reason in each case.}
\end{table}

\code{smooth\_min} (\at{wheel\_objective.py}{424}) is worth a close read. The blend
width is $k = (1 - 0.85)\cdot\max(\mathrm{cap}, 0.25)$, i.e.\ \textbf{one rung of the
exporter's fillet ladder} (the \code{max} only keeps $k$ positive when adjacent spokes
have closed the slot entirely)---narrower would model a distinction the exporter cannot
make, wider would bias the effective radius below a cap the exporter builds \emph{at}.
Exactness outside the blend is the property it exists for, and it is a claim a test
asserts. Two obvious alternatives fail it: the hyperbolic smooth-min
$\tfrac12\bigl(a + b - \sqrt{(a-b)^2 + \epsilon^2}\bigr)$ is $C^\infty$ but never exact
(measured 1.3\% below the cap at the shipped genome---inventing a fillet reduction
nothing measured), and log-sum-exp is scale-coupled and needs overflow guarding.

The docstring also records a subtle autodiff bug \emph{at exactly $a = b$}. Two
primitives are non-differentiable at the tie and autodiff picks a subgradient for each:
\code{jnp.abs} returns 0 at 0 (killing the blend term) and \code{jnp.minimum} hands the
full 1.0 to its first argument. The sum was 1.0 against a true two-sided derivative of
0.5---the function is perfectly smooth \emph{through} the tie, so this was an artifact
of the spelling, not of the mathematics. Writing the minimum symmetrically as
$\bigl(a + b - |a - b|\bigr)/2$ fixes it, confined to the blend by a \code{where}
because the symmetric form is not bit-exact. Measure-zero and never observed to bite---%
and worth fixing because the margin term drives \code{R\_hub} toward its cap on purpose,
which makes the tie an \emph{attractor} rather than an accident.

\subsection{\texttt{adjoint\_grads} --- one factorisation, \texorpdfstring{$N$}{N} solves}
\label{sec:adjointgrads}

\begin{lstlisting}[style=py]
K   = assemble_stiffness(coords, conn, u_full, ...) + contact.stiffness(coords, u_full)
Kr  = (T.T @ K @ T).tocsc()
back_solve = spla.factorized(Kr)                             # ONE LU
_, vjp = jax.vjp(lambda v: WW.mesh_coords(v, mesh), genes)   # ONE trace

for each QoI Q:
    value  = Q(c, u, y)
    dQ_dc, dQ_du, dQ_dy = jax.grad(Q, argnums=(0,1,2))(c, u, y)
    adj_r  = back_solve(-(T.T @ dQ_du))                # ONE back-substitution
    adj_f  = T @ adj_r
    d_dc, d_dy = mixed(c, u, y, adj_f, *static)        # grad_c and d/dy of adj . grad_u Pi
    grad   = vjp(dQ_dc + d_dc)[0]                      # chain to the 14 genes
\end{lstlisting}

What is shared is what is expensive: the reduced tangent $K_r$ is identical for every
quantity evaluated at the same converged state---only the right-hand side changes---and
the mesh vjp is the bigger saving at 0.774\,s versus 0.05\,s for everything else. Since
the load-control coupling needs both \code{pnorm\_stress} and \code{contact\_force} at
the same state, this is the difference between one gradient's cost and two.

There is also a guard on the adjoint solve itself: a non-finite multiplier means the
tangent at the converged state is singular, ``which means the forward solve stopped at a
state Newton could not have converged to.''

\subsection{The three searches, and the order to attack them in}
\label{sec:searches}

\code{wheel\_adjoint.py}'s docstring generalises a lesson the fillet arc paid for three
times, and it is the single most transferable idea in the repository. Three bracketed
searches appeared inside the geometry, each blocking a gradient. Each took a
\emph{different} instrument.

\begin{enumerate}
\item \textbf{\code{ring\_station}'s crossing solve}---replaced by a \emph{frozen root
  plus one Newton step}. Unrolled Newton's derivative converges to the
  implicit-function derivative, so the derivative is right without any wrapper.
\item \textbf{\code{\_fillet\_tangency}'s tangency solve}---the same treatment. A
  \code{custom\_vjp} around the bisection would have been the \emph{wrong} instrument,
  ``because there is nothing wrong with the bisection's derivative that a wrapper could
  fix. There is no derivative in a \code{sign} comparison to wrap.''
\item \textbf{The layer cliff}---90 halvings of a bracket, and the frozen-root step does
  not apply, because the answer is not a root of an equation the traced path
  re-evaluates but the \emph{output} of ninety halvings. \S85 refused the gradient
  outright. What removed the refusal was neither a wrapper nor a frozen root:
  \textbf{the cliff had a closed form and nobody had looked for one.}
\end{enumerate}

The layer width profile is a Hermite cubic that is \emph{affine in} \code{entry}:
%
\begin{equation}
  H(u) \;=\; h_{00}(u)\,w \;+\; h_{01}(u)\,(e\,w) \;+\; h_{10}(u)\,(\mathrm{entry}\cdot k)
       \;=\; a(u) \;+\; \mathrm{entry}\cdot b(u),
\end{equation}
%
writing $w$ for \code{wall}, $e$ for \code{end} and $k$ for \code{layer\_k}, with
$b(u) = u(1-u)^2 k > 0$ strictly inside $(0,1)$. So $H(u) \le Z$ exactly when
$\mathrm{entry} \le (Z - a(u))/b(u)$, and the sampled minimum reaches zero exactly when
one of them does:
%
\begin{equation}
  \mathrm{cliff} \;=\; \max_{u}\;\frac{\code{LAYER\_CLIFF\_ZERO} - a(u)}{b(u)}.
  \label{eq:cliff}
\end{equation}
%
Exact where the bisection was approximate (the two agree to \textbf{2\,ulp}, worst
relative $4\times10^{-16}$, over 98 junction pairs and 48 genomes), and differentiable
where the bisection had ninety \code{sign} comparisons. It also costs arithmetic instead
of thirty filleted sector builds, which is what makes a per-genome layer profile
adoptable at all.

\begin{quote}
\textbf{The order to try things in is therefore: is the search unnecessary; is its
answer a root worth freezing; and only then is it a discrete decision to be
classified.} Two of this project's three searches ended at the second question and one
at the first; none needed a \code{custom\_vjp}.
\end{quote}

\subsection{\texttt{hub\_fillet\_cap\_mm} --- a model built from a falsification}
\label{sec:hubcap}

The largest hub fillet the part can actually build is the \textbf{minimum of two
limits} (\at{wheel\_objective.py}{165--300}, \code{:704}), and the structure of that
\code{min} is the result of a measurement campaign (\code{make hubcap}) that overturned
the original single-term model.

\begin{itemize}
\item \textbf{The slot limit.} Two fillets grow into the gap adjacent spoke roots leave
  on the hub circle, so each may claim a share of \code{hub\_void\_deg}.
\item \textbf{The root-thickness limit.} A fillet in the re-entrant corner where a spoke
  of thickness $t_0$ meets the hub circle cannot grow past roughly half that thickness.
\end{itemize}

The 24 hub corners are \textbf{two families of twelve}, named by wedge angle rather than
by rank: a \emph{square-on} family at $266$--$270^\circ$ (the thickness mechanism) and a
\emph{near-cusp} family at $294$--$332^\circ$ (the slot mechanism). Which one binds is
\emph{observed}, not assumed---and it is not the same family on every design. At
$t_0 = 1.2$ (where everything ships) the square-on family binds; on the $t_0 = 2.55$
Stage-2 elites the near-cusp family binds.

That last fact is why the original calibration was wrong in a way its own error bars
could not show: \code{HUB\_CAP\_THICKNESS\_SHARE = 0.52} had been fitted to five ratios
inside a 0.9\% band---\textbf{all from designs where that family does not bind}. The
constant was calibrated tightly against a family that was never the limit. The result
was a cap that over-promised on all five designs by $1.07\times$ to $1.62\times$, while
the driver's own gate read green because it compared against the \emph{larger} of the
two families.

The re-fit added an arrival dependence:
%
\begin{align}
  \frac{\text{square-on threshold}}{t_0}
    &= S_{\text{thick}} \;-\; S_{\text{arr}}\,\bigl(1 - \cos\theta_{\text{hub}}\bigr)
     \nonumber\\
    &= 0.505 \;-\; 0.48\,\bigl(1 - \cos\theta_{\text{hub}}\bigr),
  \label{eq:hubcap}
\end{align}
%
writing $S_{\text{thick}}$ for \code{HUB\_CAP\_THICKNESS\_SHARE} and $S_{\text{arr}}$
for \code{HUB\_CAP\_ARRIVAL\_SLOPE}.
%
The $(1 - \cos)$ form was chosen over a marginally better quadratic \textbf{because it
stays positive across the whole physically reachable $[0^\circ, 90^\circ]$}---a
quadratic crosses zero at $87.4^\circ$ and would assert that a near-radial spoke admits
no fillet at all, which is a claim nothing measured. Out of sample on two designs in
neither fit, the law is under by 3.1--3.3\%, the same margin it carries in sample. And
the two families run in \emph{opposite} directions in arrival angle, which is why the
buildable radius is a \emph{peaked} function of it and why five designs disagreed about
the sign of the effect until a controlled sweep separated them.

The re-fit brings the five caps from $1.067 / 1.615 / 1.199 / 1.463 / 1.479$ of the
worst measured corner down to $0.979 / 0.969 / 0.827 / 0.967 / 0.969$. Note also that
\code{fillet\_cap} carries \textbf{no safety factor}, deliberately, and the comment says
why: a factor stacked on top would be a threshold with no measurement behind it. The
conservatism lives in the fit, where it is measurable.

\subsection{\texttt{t3\_terms} and the two-junction stress constraint}
\label{sec:t3}

Already covered mathematically in \S\ref{sec:stress}; two implementation notes remain.

\paragraph{The \code{(name, factory)} form is load-bearing.}
To reach \code{\_qoi\_pnorm\_stress} at $p = 4$ rather than the module default of 30, T3
passes \code{("pnorm\_stress", lambda prob: WA.\_qoi\_pnorm\_stress(prob, p=stress\_gauss\_p))}.
The bare string would take the module default and no argument could change it. The
\emph{name} stays \code{"pnorm\_stress"} so every downstream key is unchanged.

\paragraph{Probes cost nothing.}
The $p$-norm is a pure function of the converged displacement field, so any number of
exponents can be read off the field the adjoint already ran on. Each probe is pushed
through the \emph{same} \code{\_stress\_aggregate} the constraint uses, so
\code{pnorm\_by\_p[stress\_gauss\_p]} reproduces the report's own
\code{stress\_utilisation} exactly---the probe is checkable against the thing it probes.
\textbf{No gradient is taken at a probe $p$}: a convergence study reads values, and one
adjoint per exponent per phase would turn a 14-minute ladder into an hour for numbers
nothing reads.

\subsection{\texttt{selection\_key} --- which iterate of a run is promotable}
\label{sec:selectionkey}

\begin{lstlisting}[style=py]
violation = sum over BARRIER_TERMS of term value
if violation > 0:              return (2, violation, loss)
slack = hub_fillet_cap_mm - R_hub
if slack < MIN_CAP_SLACK_MM:   return (1, -slack, loss)
                               return (0, 0.0, loss)
\end{lstlisting}

Three tiers, most shippable first, sorted lexicographically. Tier 2 always ranks last
\textbf{but is still ranked}, so a run in which nothing is feasible reports its
least-violating step rather than nothing at all.

\code{MIN\_CAP\_SLACK\_MM = 1e-3} is not zero, and the reason is worth internalising:
\textbf{feasibility is fidelity-dependent.} On 2026-08-11, step 82 of a re-descent
cleared the cap at \code{medium} by 53\,nm and violated it at \code{coarse} by 93\,nm. A
barrier reading exactly 0.0 therefore certifies nothing on its own; it certifies that
\emph{this mesh} saw no violation. $10^{-3}$\,mm is four orders above that crossover and
four below what the step actually promoted was holding.

Ranking tier 1 by \emph{slack} before \emph{loss} is also deliberate: inside a
1\,$\mu$m band the loss differences are numerical noise against a feasibility question
the mesh cannot resolve, so the tiebreak that means something is distance from the knee.

\subsection{The filleted eleven-block sector}
\label{sec:filletblock}

This is the most intricate construction in the tree and it is worth understanding as a
case study in structured meshing, even though it is not yet wired into the optimizer.

\paragraph{The problem.}
Rounding the $P_t$ corner \emph{inside} the spoke block (the original approach, retired
by \S47) does not remove the fold---it moves it into whichever block carries the arc,
where it becomes a cusp. The fix re-cuts the whole sector.

\paragraph{The construction.}
Per junction, two new blocks:

\begin{itemize}
\item \code{<j>\_fillet\_a}---a boundary layer along the fillet arc, above the ring
  circle;
\item \code{<j>\_fillet\_b}---the wedge that crosses the ring circle, out to the ring's
  \emph{far} side.
\end{itemize}

Plus the junction block (unchanged in shape, but its end cross-section is replaced by
\code{fillet\_a}'s inner edge) and the two ring blocks, re-cut. Seven blocks becomes
eleven, fourteen whole-edge seams, closing at $3\times10^{-14}$\,mm.

\paragraph{Why two fillet blocks and not one.}
$N$ is where the fillet block's inner edge crosses the ring circle. Above $N$ that
edge's partner is the junction block; below it, the ring. \textbf{One edge with two
partners is a partial-edge seam}, and whole-edge single ownership is the module's whole
safety net---so the block is split instead.

\paragraph{Why the cut must reach the ring's far boundary.}
This is a nice piece of reasoning. A cut stopping partway gives the ring free block's
left edge two partners. Splitting \emph{that} block gives its own right edge two
partners, so the split propagates round the ring. And the block it propagates into is a
\textbf{triangle}, because the inner edge is a concentric offset of an arc tangent to
the ring circle, and is therefore tangent to every concentric circle ($12.86^\circ$ at
the hub, $4.38^\circ$ at the rim, measured). Carrying the cut all the way through splits
the ring into two quadrilaterals and \emph{terminates}. The price: the ring's radial
node count becomes \code{n\_thick}, so \code{n\_collar\_r} and \code{n\_rim\_r} are
unused by this blocking.

\paragraph{The \code{entry} slope is negative on purpose.}
At $\code{entry} = 0$ the inner edge leaves the end cross-section \emph{tangent} to the
far flank---which is the junction block's own top edge---and three blocks meet at that
node with $180^\circ$ to share between two of them. Measured, the junction block's
minimum scaled Jacobian is 0.0400 there against 0.4272 at the chosen slope.

\paragraph{Where it stands.}
As of \S89--\S92, no \emph{mesh-validity} objection to wiring it in survives: the
filleted and unfilleted meshes score the same 31-of-32 on a held-out draw against
\code{MIN\_SJ\_TARGET}, and the one genome under the barrier is under it \emph{harder}
unfilleted (0.1298 versus 0.1775; barrier term 327.17 versus 18.20). What the fillet
costs is \textbf{headroom, not crossings}---median minimum scaled Jacobian 0.338 against
0.745. Measured costs: $1.12\times$ per objective evaluation (not the long-quoted
2--3$\times$), with the one-time JIT trace at $4.13\times$ and nineteen minutes. The two
meshes disagree about the solved wheel by 37.97\% of axle drop (linear, one phase) and
47.85\% (SVK, eight phases).
\code{tests/test\_corner\_singularity.py::}\allowbreak
\code{test\_nothing\_wires\_the\_fillet\_into\_the\_objective}
holds the scope gate as a \emph{check}, so it is a deliberate decision rather than an
accident.

% ===========================================================================
\section{How it all fits together}
\label{sec:synthesis}
% ===========================================================================

\subsection{One idea, applied at every level: write the thing once, derive the rest}

The repository has a single organizing principle and applies it recursively.

\begin{itemize}
\item \textbf{In the element:} write the strain energy. The internal force is its
  gradient and the tangent is its Hessian, both by \code{jax.grad}/\code{jax.hessian}.
  No hand-derived tangent can drift.
\item \textbf{In the contact law:} write the penalty potential. The contact force, the
  contact stiffness, \emph{and the total ground reaction} are all derivatives of it. The
  reaction agreeing with an independent pressure quadrature is then evidence, not a
  tautology.
\item \textbf{In the sensitivity:} the residual is already $\nabla\Pi$, so the adjoint
  is one \code{jax.grad} of one scalar. No separately coded sensitivity operator.
\item \textbf{In the geometry:} \code{offset\_band} is written once and serves both the
  exported outline ($\code{n\_across} = 1$) and the mesh block
  ($\code{n\_across} = 6$), so the meshed part and the exported part are the same
  surface by construction.
\item \textbf{In the constants:} \code{MAX\_ARRIVAL\_DEG}, \code{MIN\_FOLD\_MARGIN\_MM},
  \code{TAPER\_BREAKPOINTS} and \code{MIN\_JUNCTION\_BITE} each live in exactly one
  module, chosen so that both sides of the interpreter boundary can import it.
\end{itemize}

The consequences compound. Because the energy is written once, the linear and SVK
kinematics share a material law and cannot disagree about it. Because the residual
\emph{is} a gradient, the Newton line search's slope is exact and a non-negative value
diagnoses instability rather than noise. Because the geometry is closed-form and the
connectivity is a compile-time constant, the whole pipeline from 14 genes to an axle
drop is differentiable---which is what makes Stage 3 possible at all.

\subsection{The second idea: name what is not differentiable}

The tree is unusually disciplined about the boundary between smooth and discrete. Four
things are frozen, and every one is \emph{named} rather than tolerated:

\begin{enumerate}
\item \textbf{Flank orientation}---a topological fact about the genome; a step that
  flips it is a real discontinuity of the design space.
\item \textbf{Seam ownership}---a labelling, verified by \code{check\_seams} to have
  made no difference.
\item \textbf{Fillet roots and refusals}---frozen from the eager build, refined by one
  Newton step.
\item \textbf{\code{buckling}}---the legacy Euler proxy, with an \emph{asserted} zero
  gradient and a gate that fails the day it stops being zero.
\end{enumerate}

And two things are named as \emph{absent}: \code{R\_hub} and \code{R\_rim} have
identically zero FEA gradient on the unfilleted mesh, because no fillet is meshed---so
$\partial\,\mathrm{coords}/\partial\,\mathrm{gene}$ is zero before any solver is
involved. This is not left to be discovered; \code{insensitive\_genes}
(\at{wheel\_adjoint.py}{744}) measures it on the Jacobian with $\mathrm{tol} = 0.0$
(``these columns are not small, they are absent, and a tolerance would invite them to
become small''), and the census is gated. Those two genes still carry a live gradient
through the closed-form \code{fillet} and \code{fillet\_cap} terms, which is the whole
reason those terms exist.

\subsection{The third idea: a measurement is not a number, it is a number plus its scope}

Almost every constant in this tree carries the measurement that produced it, the sample
it was measured on, and---repeatedly---the \emph{previous} value and why it was wrong.
The pattern that recurs:

\begin{itemize}
\item \code{HUB\_CAP\_THICKNESS\_SHARE} was fitted tightly to a corner family that never
  binds.
\item \code{MIN\_FOLD\_MARGIN\_MM} is 0.1 and not 0 because zero admits six inverted
  meshes.
\item \code{T1\_REJECT} is $10^4$ and not 1.0 because 1.0 deadlocked a probe for
  thirty-five minutes.
\item \code{EMBED\_ALLOWANCE\_PER\_SPOKE\_MM2 = 3.03} is \textbf{deliberately not
  updated}, because the same computation has read 4.356, 3.032, 0.98 and 0.5317 on four
  genomes---``what is needed is the scaling law'', and replacing the constant would only
  re-stale it on the next genome.
\item \code{FILLET\_LAYER\_CLIFF\_FACTOR = 0.45} is the \emph{bottom edge} of an
  admissible band, because two axes improve monotonically across a flat third---and
  quoting the interior of a band is exactly the error a previous section made.
\end{itemize}

The \code{xfail\_strict = true} setting in \code{pyproject.toml} is the same idea
mechanised: an accepted finding is recorded as a strict \code{xfail} naming the section
that decided it, so \textbf{the day the wheel changes enough for it to pass, the suite
reopens the question by itself}.

\subsection{Reading the tree}

If you want to work in this repository, here is the order I would read in.

\begin{enumerate}
\item \code{CLAUDE.md}, then \code{PLAN.md}'s header block and the ``Open arcs''
  table---that is the ranking and the rules.
\item \code{src/wheel\_geometry.py} end to end. It is 436 lines and it is the vocabulary
  for everything else: B\'ezier, taper, offset band, fold margin, arrival.
\item \code{src/wheel\_wheel.py}'s module docstring (lines 1--170). It explains why a
  structured partition of this wheel exists at all, what is and is not modelled, and
  what a seam mismatch would silently do.
\item \code{src/wheel\_fem.py}'s module docstring, then \code{element\_energy} and
  \code{\_strain}. Sixty lines gets you the whole physics.
\item \code{src/wheel\_adjoint.py}'s module docstring---the adjoint derivation and the
  ``unnecessary / freezable / discrete'' ordering for searches.
\item \code{src/wheel\_objective.py}'s module docstring for the T1/T2/T3 economics, then
  jump to the term you care about.
\item Whichever \code{studies/study\_*.py} docstring names the number you are about to
  quote. Every one of them says what it measured, on what, and what it does \emph{not}
  establish.
\end{enumerate}

And one habit the tree will reward: \textbf{before quoting any per-design number, check
which genome it was measured on.} The provenance chain is \code{36aed36} (the GA/beam
optimum, pinned forever) $\to$ \code{350f4c7} $\to$ \code{e4219f3} $\to$ \code{e126cc3}
$\to$ \code{09e8188} (shipped). Several docstrings were written under earlier links and
say so at the top. The banner in \code{PLAN.md} went stale for two promotions, and that
fact is itself recorded as the warning.

\vfill
\begin{center}\rule{0.45\linewidth}{0.4pt}\end{center}
\begin{center}\small\itshape
Written from the source. Line references are against the \code{feature} branch as of
2026-08-29; \code{PLAN.md} sections referenced up to \S92.
\end{center}

\end{document}
